\documentclass[12pt,a4paper]{article}
\usepackage[utf8]{inputenc}
\usepackage[T5]{fontenc}
\usepackage[vietnamese]{babel}
\usepackage{amsmath, amssymb}
\usepackage{array}
\usepackage{enumitem}
\usepackage{setspace}
\setstretch{1.18}
\usepackage{graphicx}
\usepackage{ifthen}
\usepackage{tikz}
\usetikzlibrary{arrows.meta, positioning, patterns}
\IfFileExists{pgf-pie.sty}{\usepackage{pgf-pie}}{}
\usepackage{pgfplots}
\pgfplotsset{compat=1.18}
\usepackage{chemfig}
\usepackage{longtable,booktabs}
\usepackage{adjustbox}
\providecommand{\errorpart}[2]{$\underset{\textbf{#1}}{\underline{\text{#2}}}$}
\usepackage[top=0.8in, left=0.75in, right=0.75in, bottom=0.85in, headheight=38pt, headsep=18pt, footskip=1.2cm]{geometry}
\usepackage{fancyhdr}
\usepackage[colorlinks=true, linkcolor=blue, urlcolor=blue]{hyperref}

\setlength{\parindent}{0pt}
\setlength{\parskip}{5pt plus 2pt minus 1pt}

\setlist[enumerate]{itemsep=4pt, parsep=2pt, topsep=4pt, partopsep=2pt}
\setlist[itemize]{itemsep=4pt, parsep=2pt, topsep=4pt}

\makeatletter
\AtEndDocument{\immediate\write\@auxout{\string\newlabel{LecturingLastPage}{{}{\thepage}{}{}{}}}}
\AtEndDocument{\immediate\write\@auxout{\string\newlabel{LecturingPageCount}{{}{\two@digits{\value{page}}}{}{}{}}}}
\makeatother
\pagestyle{fancy}
\fancyhf{}
\cfoot{\small \thepage}
\renewcommand{\headrulewidth}{0pt}
\renewcommand{\footrulewidth}{0pt}

\begin{document}

\sloppy

\begin{center}
{\bfseries \mbox{}}\\[0.1cm]
{\bfseries KỲ THI ĐÁNH GIÁ NĂNG LỰC}\\[0.1cm]
{\bfseries ĐỀ LUYỆN TẬP}
\end{center}
\hfill{\small Mã đề: 9245}

\noindent\begin{tabular}{@{}p{0.28\textwidth}p{0.68\textwidth}@{}}
\textbf{Họ và tên thí sinh:} & \dotfill \\
\textbf{Số báo danh:} & \dotfill \\
\textbf{Thời gian làm bài:} & 150 phút (không kể thời gian phát đề) \\
\textbf{Tổng số câu hỏi:} & 120 câu hỏi \\
\textbf{Dạng câu hỏi:} & Trắc nghiệm, với 04 phương án lựa chọn (Trong đó, chỉ có 01 phương án đúng) \\
\textbf{Cách làm bài:} & Tô đậm phương án đã chọn vào phiếu trả lời trắc nghiệm \\
\end{tabular}

\begin{center}\textbf{CẤU TRÚC ĐỀ THI}\\[0.1cm]
\begin{tabular}{|p{0.5\textwidth}|c|c|}\hline
\textbf{Nội dung} & \textbf{Số câu} & \textbf{Thứ tự câu} \\ \hline
\textbf{Phần 1: Sử dụng ngôn ngữ} & 60 & 1–60 \\ \hline
\hspace{0.5cm}1.1. Tiếng Việt & 30 & 1–30 \\ \hline
\hspace{0.5cm}1.2. Tiếng Anh & 30 & 31–60 \\ \hline
\textbf{Phần 2: Toán học} & 30 & 61–90 \\ \hline
\textbf{Phần 3: Tư duy khoa học} & 30 & 91–120 \\ \hline
\hspace{0.5cm}3.1. Logic, phân tích số liệu & 12 & 91–102 \\ \hline
\hspace{0.5cm}3.2. Suy luận khoa học & 18 & 103–120 \\ \hline
\end{tabular}\end{center}
\vspace{0.4cm}

% LECTURING-PART: TV
\textbf{PHẦN 1: SỬ DỤNG NGÔN NGỮ\\ 1.1. TIẾNG VIỆT} \textit{}

\makeatletter\@beginparpenalty=10000 \@endparpenalty=10000 \makeatother

% LECTURING-BLOCK: single

\interlinepenalty=10000

\noindent \textbf{Câu 1:} Trong bài thơ \textit{Tự tình II}, Hồ Xuân Hương viết:
\begin{center}
``Xiên ngang mặt đất, rêu từng đám,\\*
Đâm toạc chân mây, đá mấy hòn.''
\end{center}
\begin{flushright}\small\textit{(Hồ Xuân Hương, Tự tình II)}\end{flushright}
Cách dùng các từ ``xiên ngang'', ``đâm toạc'' trong hai câu thơ có tác dụng gì?\\
\noindent A. Tả cảnh thiên nhiên tĩnh lặng để làm nền cho nỗi buồn cô đơn của nhân vật trữ tình.\\
\noindent B. Thể hiện sự hài hòa, tròn đầy giữa con người với thiên nhiên.\\
\noindent C. Nhấn mạnh sự nhỏ bé, cam chịu của rêu và đá trước đất trời.\\
\noindent D. Gợi sức sống mãnh liệt của cảnh vật, đồng thời bộc lộ thái độ bứt phá, phản kháng của nhân vật trữ tình trước số phận.

\interlinepenalty=0


\vspace{0.3cm}
\interlinepenalty=10000

\noindent \textbf{Câu 2:} Đọc bài thơ sau:
\begin{center}
Ao thu lạnh lẽo nước trong veo,\\*
Một chiếc thuyền câu bé tẻo teo.\\*
Sóng biếc theo làn hơi gợn tí,\\*
Lá vàng trước gió khẽ đưa vèo.\\*
Tầng mây lơ lửng trời xanh ngắt,\\*
Ngõ trúc quanh co khách vắng teo.\\*
Tựa gối buông cần lâu chẳng được,\\*
Cá đâu đớp động dưới chân bèo.
\end{center}
\begin{flushright}\small\textit{(Nguyễn Khuyến, Thu điếu)}\end{flushright}
Nhận xét nào đúng về thể thơ và cách gieo vần của bài thơ?\\
\noindent A. Thất ngôn bát cú, vần lưng ``eo'' gieo ở tiếng thứ tư của mỗi câu.\\
\noindent B. Thất ngôn bát cú, vần chân ``eo'' gieo ở cuối các câu 2, 4, 6, 8.\\
\noindent C. Thất ngôn bát cú, vần chân ``eo'' gieo ở cuối các câu 1, 2, 4, 6, 8.\\
\noindent D. Thất ngôn trường thiên, vần ``eo'' gieo liên tục ở mọi câu.

\interlinepenalty=0


\vspace{0.3cm}
\interlinepenalty=10000

\noindent \textbf{Câu 3:} Đọc đoạn thơ sau:
\begin{center}
Đêm thu buồn lắm chị Hằng ơi!\\*
Trần thế em nay chán nửa rồi,\\*
Cung quế đã ai ngồi đó chửa?\\*
Cành đa xin chị nhắc lên chơi.
\end{center}
\begin{flushright}\small\textit{(Tản Đà, Muốn làm thằng Cuội)}\end{flushright}
Nhận xét nào đúng về cách xưng hô và giọng điệu của đoạn thơ?\\
\noindent A. Cách xưng ``em'' – gọi ``chị Hằng'' tạo giọng thân mật, hóm hỉnh, bộc lộ cái tôi lãng mạn chán cõi trần và khao khát thoát li.\\
\noindent B. Giọng hào sảng của người anh hùng khẳng định chí lớn giữa đời.\\
\noindent C. Cách xưng hô cung kính, e dè của người phàm trước thần linh, thể hiện lòng ngưỡng vọng cõi tiên.\\
\noindent D. Giọng ai oán, nghẹn ngào của người chinh phụ than thân trách phận.

\interlinepenalty=0


\vspace{0.3cm}
\interlinepenalty=10000

\noindent \textbf{Câu 4:} Trong \textit{Văn tế nghĩa sĩ Cần Giuộc}, Nguyễn Đình Chiểu viết về những người nông dân nghĩa sĩ:
\begin{center}
``Ngoài cật có một manh áo vải, nào đợi mang bao tấu, bầu ngòi,\\*
Trong tay cầm một ngọn tầm vông, chỉ nài sắm dao tu, nón gõ.''
\end{center}
\begin{flushright}\small\textit{(Nguyễn Đình Chiểu, Văn tế nghĩa sĩ Cần Giuộc)}\end{flushright}
Hai câu văn có đặc điểm nghệ thuật và ý nghĩa nào sau đây?\\
\noindent A. Dùng câu văn biền ngẫu đối nhau, gợi sự thiếu thốn trang bị mà lòng tự nguyện, ý chí chiến đấu vẫn sục sôi.\\
\noindent B. Dùng lối liệt kê để miêu tả quá trình luyện tập bài bản của nghĩa quân trước trận đánh.\\
\noindent C. Dùng hình ảnh tượng trưng để ca ngợi tài thao lược của người chỉ huy nghĩa quân.\\
\noindent D. Dùng lối nói phóng đại để giễu nhại sự non nớt của nghĩa quân trong chiến trận.

\interlinepenalty=0


\vspace{0.3cm}
\interlinepenalty=10000

\noindent \textbf{Câu 5:} Khi bị giam lỏng ở lầu Ngưng Bích, Thúy Kiều nhìn ra xa:
\begin{center}
``Buồn trông cửa bể chiều hôm,\\*
Thuyền ai thấp thoáng cánh buồm xa xa?\\*
Buồn trông ngọn nước mới sa,\\*
Hoa trôi man mác, biết là về đâu?''
\end{center}
\begin{flushright}\small\textit{(Nguyễn Du, Truyện Kiều)}\end{flushright}
Việc lặp lại cụm từ ``buồn trông'' ở đầu các câu lục có tác dụng chủ yếu nào sau đây?\\
\noindent A. Liệt kê các cảnh vật khác nhau để giới thiệu khách quan phong cảnh vùng cửa bể.\\
\noindent B. Tạo nhịp điệu trùng điệp như những đợt sóng lòng, nhấn mạnh nỗi buồn tăng dần và thân phận trôi nổi, vô định của Kiều.\\
\noindent C. Biểu lộ thái độ thờ ơ của Kiều trước cảnh vật xung quanh.\\
\noindent D. Chuyển mạch kể từ cảnh thực sang một giấc mơ hoang đường.

\interlinepenalty=0


\vspace{0.3cm}
\interlinepenalty=10000

\noindent \textbf{Câu 6:} Nhân vật ``tôi'' (ông giáo) trong truyện ngắn \textit{Lão Hạc} suy ngẫm:\\
``Chao ôi! Đối với những người ở quanh ta, nếu ta không cố tìm mà hiểu họ, thì ta chỉ thấy họ gàn dở, ngu ngốc, bần tiện, xấu xa, bỉ ổi... toàn những cớ để cho ta tàn nhẫn; không bao giờ ta thấy họ là những người đáng thương: không bao giờ ta thương...''
\begin{flushright}\small\textit{(Nam Cao, Lão Hạc)}\end{flushright}
Đoạn văn thể hiện điều gì?\\
\noindent A. Lời tự trách của ông giáo vì đã đối xử tàn nhẫn với Lão Hạc.\\
\noindent B. Lời phê phán những người nghèo vì tính gàn dở, bần tiện của họ.\\
\noindent C. Lời khẳng định rằng con người vốn xấu xa và không đáng được thương.\\
\noindent D. Lời triết lí về việc phải cố gắng thấu hiểu hoàn cảnh con người thì mới có thể cảm thông và yêu thương họ.

\interlinepenalty=0


\vspace{0.3cm}
\interlinepenalty=10000

\noindent \textbf{Câu 7:} Đọc đoạn văn sau:\\
``Tiếng trống thu không trên cái chợ của huyện nhỏ; từng tiếng một vang xa để gọi buổi chiều. Phương tây đỏ rực như lửa cháy và những đám mây ánh hồng như hòn than sắp tàn. Dãy tre làng trước mặt đen lại và cắt hình rõ rệt trên nền trời.''
\begin{flushright}\small\textit{(Thạch Lam, Hai đứa trẻ)}\end{flushright}
Hình ảnh ``những đám mây ánh hồng như hòn than sắp tàn'' có tác dụng gì?\\
\noindent A. Dùng phép nhân hóa làm đám mây trở nên có tâm trạng như con người.\\
\noindent B. Dùng phép liệt kê để tả đầy đủ các màu sắc của bầu trời ban ngày.\\
\noindent C. Dùng phép ẩn dụ để nói về sự ấm áp, sum vầy của gia đình.\\
\noindent D. Dùng phép so sánh gợi vẻ rực rỡ mà sắp tàn của chiều quê, báo trước nỗi buồn man mác của nhân vật.

\interlinepenalty=0


\vspace{0.3cm}
\interlinepenalty=10000

\noindent \textbf{Câu 8:} Trong \textit{Bình Ngô đại cáo} (bản dịch của Ngô Tất Tố) có đoạn:
\begin{center}
``Như nước Đại Việt ta từ trước,\\*
Vốn xưng nền văn hiến đã lâu.\\*
Núi sông bờ cõi đã chia,\\*
Phong tục Bắc Nam cũng khác.\\*
Từ Triệu, Đinh, Lý, Trần bao đời xây nền độc lập,\\*
Cùng Hán, Đường, Tống, Nguyên mỗi bên hùng cứ một phương.\\*
Tuy mạnh yếu có lúc khác nhau,\\*
Song hào kiệt thời nào cũng có.''
\end{center}
\begin{flushright}\small\textit{(Nguyễn Trãi, Bình Ngô đại cáo)}\end{flushright}
Căn cứ nào sau đây \textbf{không} được nêu trong đoạn trên để khẳng định Đại Việt là một quốc gia độc lập?\\
\noindent\makebox[0.48\textwidth][l]{A. Cương vực lãnh thổ riêng.}\makebox[0.48\textwidth][l]{B. Nền văn hiến lâu đời.}\\
\noindent\makebox[0.48\textwidth][l]{C. Ngôn ngữ và chữ viết riêng.}\makebox[0.48\textwidth][l]{D. Phong tục tập quán riêng.}

\interlinepenalty=0


\vspace{0.3cm}
\interlinepenalty=10000

\noindent \textbf{Câu 9:} Bài \textit{Văn tế nghĩa sĩ Cần Giuộc} mở đầu như sau:
\begin{center}
``Hỡi ơi!\\*
Súng giặc đất rền;\\*
Lòng dân trời tỏ.\\*
Mười năm công vỡ ruộng, chưa ắt còn danh nổi như phao;\\*
Một trận nghĩa đánh Tây, tuy là mất tiếng vang như mõ.''
\end{center}
\begin{flushright}\small\textit{(Nguyễn Đình Chiểu, Văn tế nghĩa sĩ Cần Giuộc)}\end{flushright}
Đoạn trên thuộc phần nào của bài văn tế và có nội dung chính là gì?\\
\noindent A. Phần thích thực, kể lại cuộc đời và công trạng của những nghĩa sĩ.\\
\noindent B. Phần lung khởi, nêu ý nghĩa lớn lao của cái chết vì nghĩa đối với lòng dân.\\
\noindent C. Phần kết, cầu mong linh hồn nghĩa sĩ được yên nghỉ.\\
\noindent D. Phần ai vãn, bày tỏ nỗi xót thương của người đứng tế.

\interlinepenalty=0


\vspace{0.3cm}
\interlinepenalty=10000

\noindent \textbf{Câu 10:} Đọc đoạn thơ sau:
\begin{center}
``Trăm năm trong cõi người ta,\\*
Chữ tài chữ mệnh khéo là ghét nhau.\\*
Trải qua một cuộc bể dâu,\\*
Những điều trông thấy mà đau đớn lòng.''
\end{center}
\begin{flushright}\small\textit{(Nguyễn Du, Truyện Kiều)}\end{flushright}
Nhận định nào sau đây phù hợp với nội dung đoạn thơ?\\
\noindent A. Tài năng giúp con người vượt qua mọi biến cố; ``bể dâu'' chỉ vùng đất nuôi dưỡng nhân tài.\\
\noindent B. Số phận do trời định sẵn nên con người không cần rèn luyện tài năng; ``bể dâu'' chỉ nghề trồng dâu nuôi tằm.\\
\noindent C. Tài năng thường đi kèm số phận trắc trở; ``bể dâu'' chỉ những biến đổi lớn lao của cuộc đời khiến người trải nghiệm xót xa.\\
\noindent D. Cuộc đời an nhàn là phần thưởng của người có tài; ``bể dâu'' chỉ cảnh thái bình thịnh vượng.

\interlinepenalty=0


\vspace{0.3cm}
\interlinepenalty=10000

\noindent \textbf{Câu 11:} Dãy từ nào dưới đây viết đúng chính tả?\\
\noindent\makebox[0.48\textwidth][l]{A. Chặt chẽ, trau chuốt, xuất xắc.}\makebox[0.48\textwidth][l]{B. Chặt chẽ, trau chuốt, xuất sắc.}\\
\noindent\makebox[0.48\textwidth][l]{C. Chặt chẽ, chau chuốt, suất sắc.}\makebox[0.48\textwidth][l]{D. Chặc chẽ, trau chuốt, xuất sắc.}

\interlinepenalty=0


\vspace{0.3cm}
\interlinepenalty=10000

\noindent \textbf{Câu 12:} Câu nào sau đây là câu đúng ngữ pháp?\\
\noindent A. Với tinh thần trách nhiệm cao, đã giúp anh hoàn thành tốt nhiệm vụ được giao.\\
\noindent B. Qua tác phẩm ``Chí Phèo'' cho thấy sự tha hóa của người nông dân trước Cách mạng.\\
\noindent C. Hội nghị đã thảo luận và thông qua kế hoạch hoạt động của năm học mới.\\
\noindent D. Người con gái ấy có đôi mắt to, hiền lành và rất chăm chỉ.

\interlinepenalty=0


\vspace{0.3cm}
\interlinepenalty=10000

\noindent \textbf{Câu 13:} Trong câu ``Anh ấy vẫn giữ thái độ bàng quan trước những bất công ngay trong cơ quan'', từ ``bàng quan'' có nghĩa là gì?\\
\noindent A. Đứng ngoài cuộc, thờ ơ, không quan tâm đến việc đang xảy ra.\\
\noindent B. Lo lắng, bất an trước những điều chưa rõ.\\
\noindent C. Quan sát sự việc từ nhiều phía để đưa ra phán xét.\\
\noindent D. Nhìn nhận sự việc một cách khách quan, công bằng.

\interlinepenalty=0


\vspace{0.3cm}
\interlinepenalty=10000

\noindent \textbf{Câu 14:} Chọn cặp thuật ngữ phù hợp để điền vào chỗ trống: ``Trong thơ Đường luật, \underline{\hspace{1cm}} là quy tắc làm cho các câu thơ dính nhau theo thanh bằng – trắc, còn \underline{\hspace{1cm}} là sự hòa âm bằng những tiếng có phần vần giống nhau.''\\
\noindent\makebox[0.24\textwidth][l]{A. niêm – nhịp}\makebox[0.24\textwidth][l]{B. đối – vần}\makebox[0.24\textwidth][l]{C. niêm – vần}\makebox[0.24\textwidth][l]{D. luật – đối}

\interlinepenalty=0


\vspace{0.3cm}
\interlinepenalty=10000

\noindent \textbf{Câu 15:} Thành ngữ nào phù hợp để điền vào chỗ trống trong câu: ``Cả ngày anh ta \underline{\hspace{1.5cm}}, chẳng chịu làm nên trò trống gì.''?\\
\noindent\makebox[0.48\textwidth][l]{A. ăn bữa nay lo bữa mai}\makebox[0.48\textwidth][l]{B. ăn xổi ở thì}\\
\noindent\makebox[0.48\textwidth][l]{C. ăn không ngồi rồi}\makebox[0.48\textwidth][l]{D. ăn cháo đá bát}

\interlinepenalty=0


\vspace{0.3cm}
\interlinepenalty=10000

\noindent \textbf{Câu 16:} Câu nào sau đây thể hiện phong cách ngôn ngữ hành chính?\\
\noindent A. Căn cứ Luật Giáo dục năm 2019, Hiệu trưởng quyết định thành lập Hội đồng coi thi kỳ thi học kỳ I.\\
\noindent B. Hôm qua tui nghe nói thầy cho nghỉ sớm rồi đó, mai đi học không?\\
\noindent C. Giọt sương mai long lanh như hạt ngọc nằm yên trên cành lá.\\
\noindent D. Nước tinh khiết sôi ở 100°C khi áp suất khí quyển là 1 atm.

\interlinepenalty=0


\vspace{0.3cm}
\interlinepenalty=10000

\noindent \textbf{Câu 17:} Sau cái chết của Lão Hạc, nhân vật ``tôi'' nói:\\
``Không! Cuộc đời chưa hẳn đã đáng buồn, hay vẫn đáng buồn nhưng lại đáng buồn theo một nghĩa khác.''
\begin{flushright}\small\textit{(Nam Cao, Lão Hạc)}\end{flushright}
Câu văn thể hiện sự thay đổi nhận thức nào của nhân vật ``tôi''?\\
\noindent A. Từ chỗ buồn vì mất một người bạn thân, chuyển sang vui vì Lão Hạc đã được giải thoát.\\
\noindent B. Từ chỗ tin vào lòng tốt của con người, chuyển sang hoàn toàn tuyệt vọng về cuộc đời.\\
\noindent C. Từ chỗ trách Lão Hạc không hiểu mình, chuyển sang trách bản thân đã không giúp đỡ lão.\\
\noindent D. Từ chỗ buồn vì thấy con người có thể tha hóa, chuyển sang nhận ra nỗi buồn thật sự là người tốt bị cái đói nghèo dồn đến đường cùng.

\interlinepenalty=0


\vspace{0.3cm}
\interlinepenalty=10000

\noindent \textbf{Câu 18:} Chọn từ thích hợp điền vào chỗ trống: ``Bản báo cáo được \underline{\hspace{1.5cm}} kĩ lưỡng về số liệu trước khi trình lên lãnh đạo.''\\
\noindent\makebox[0.24\textwidth][l]{A. thẩm vấn}\makebox[0.24\textwidth][l]{B. thẩm thấu}\makebox[0.24\textwidth][l]{C. thẩm định}\makebox[0.24\textwidth][l]{D. thẩm mĩ}

\interlinepenalty=0


\vspace{0.3cm}
\interlinepenalty=10000

\noindent \textbf{Câu 19:} Câu nào sau đây biểu thị quan hệ tương phản giữa hai vế?\\
\noindent A. Anh viết báo cáo để kịp nộp vào sáng mai.\\
\noindent B. Trời đã khuya nên anh dừng viết báo cáo.\\
\noindent C. Nếu trời khuya thì anh sẽ dừng viết báo cáo.\\
\noindent D. Trời đã khuya mà anh vẫn miệt mài viết báo cáo.

\interlinepenalty=0


\vspace{0.3cm}
\interlinepenalty=10000

\noindent \textbf{Câu 20:} Từ ``đi'' trong câu nào dưới đây được dùng theo cách nói giảm, nói tránh để chỉ cái chết?\\
\noindent A. Chiếc đồng hồ này đi chậm hơn đồng hồ chuẩn năm phút.\\
\noindent B. Em bé đang chập chững tập đi trong sân nhà.\\
\noindent C. Ông nội tôi đã đi từ mùa thu năm ngoái, để lại khu vườn đầy hoa cúc.\\
\noindent D. Bà đi chợ từ sáng sớm để mua cá.

\interlinepenalty=0


\vspace{0.5cm}
% LECTURING-BLOCK: passage
\noindent\textbf{Dựa vào thông tin dưới đây để trả lời các câu từ 21 đến 25}% LECTURING-BLOCK-INSTRUCTION


\par\nobreak

(1) Nét độc đáo và đặc điểm chính của chợ nổi Cái Răng là chuyên mua bán các loại rau củ, trái cây của vùng đồng bằng sông Cửu Long. Thuở xưa, chợ nổi hình thành vì hệ thống giao thông đường thủy giữ vai trò gần như độc tôn. Trong khi nhu cầu thương mại ngày càng tăng cao mà việc đi lại luôn gắn với dòng sông, bến nước nên điểm giao sông, bùng binh trở thành địa điểm lý tưởng để người ta tụ tập mua bán. Ngày nay, hệ thống giao thông đường bộ đã phát triển rộng khắp, điều này góp phần làm ``biến mất'' nhiều chợ nổi. Tuy nhiên, dưới tác động của hoạt động du lịch, chợ nổi dần tách khỏi vai trò chính yếu của nó để tồn tại và phát triển theo một hướng mới bền vững hơn.

(2) Chợ nổi Cái Răng cũng là chợ đầu mối chuyên mua bán các loại rau củ của vùng. Chợ nhóm từ 2 - 3 giờ sáng, nhộn nhịp nhất vào khoảng 4 - 6 giờ sáng, hoạt động mua bán diễn ra cả ngày. Vì là chợ đầu mối nên hàng hóa tập trung ở đây với số lượng tương đối lớn. Mỗi mặt hàng đã được phân loại cho đồng đều về chất lượng, kích cỡ. Phương tiện di chuyển và chuyên chở hàng hóa chủ yếu là ghe bầu.

(3) Hòa mình vào không khí nhộn nhịp của buổi chợ, du khách có thể quan sát, tìm hiểu sinh hoạt của nhiều gia đình thương hồ với nhiều thế hệ chung sống trên ghe. Mỗi chiếc ghe là một căn nhà trên sông nước với những chậu hoa kiểng, các loài vật nuôi, các tiện nghi đầy đủ như ti-vi màu, đầu dĩa, dàn âm thanh...có cả xe gắn máy đậu trên ghe.

(4) Hình thức mua bán trên chợ là treo ``bẹo'' thay vì treo biển hiệu. Thông thường, ghe bán sẽ dựng một chậu cây bất kỳ để phân biệt với ghe mua.

(5) Chợ nổi Cái Răng là một trong những điểm tham quan đặc sắc nhất ở Cần Thơ. Có thể xem chợ nổi là bản sắc văn hóa của vùng đồng bằng sông nước, thu hút rất nhiều du khách, đặc biệt là du khách nước ngoài. Hầu như các tour tham quan TP. Cần Thơ đều có điểm đến là chợ nổi Cái Răng.
\begin{flushright}\small\textit{(Theo Wikipedia tiếng Việt, mục ``Chợ nổi Cái Răng'', vi.wikipedia.org, truy cập 08/10/2026; có lược bớt một số mục)}\end{flushright}

\par\nobreak

\interlinepenalty=10000

\noindent \textbf{Câu 21:} Theo đoạn (1), vì sao chợ nổi hình thành từ xưa?\\
\noindent A. Vì giao thông đường thủy gần như độc tôn nên điểm giao sông trở thành nơi lý tưởng để tụ tập mua bán.\\
\noindent B. Vì hoạt động du lịch thu hút khách nước ngoài đến mua trái cây.\\
\noindent C. Vì các gia đình thương hồ muốn có nơi sinh sống cố định trên sông.\\
\noindent D. Vì giao thông đường bộ phát triển nên việc buôn bán dời xuống sông.

\interlinepenalty=0


\vspace{0.3cm}
\interlinepenalty=10000

\noindent \textbf{Câu 22:} Thông tin nào ở đoạn (2) cho thấy chợ nổi Cái Răng mang tính chất của một ``chợ đầu mối''?\\
\noindent A. Hàng hóa tập trung với số lượng tương đối lớn và đã được phân loại đồng đều về chất lượng, kích cỡ.\\
\noindent B. Chợ họp từ rất sớm và hoạt động trong cả ngày.\\
\noindent C. Phương tiện chuyên chở chủ yếu là ghe bầu.\\
\noindent D. Chợ chuyên bán các loại rau củ của vùng.

\interlinepenalty=0


\vspace{0.3cm}
\interlinepenalty=10000

\noindent \textbf{Câu 23:} Việc liệt kê ``ti-vi màu, đầu dĩa, dàn âm thanh'', ``xe gắn máy'' ở đoạn (3) nhằm mục đích gì?\\
\noindent A. Cho thấy gia đình thương hồ sống đầy đủ tiện nghi ngay trên chiếc ghe, như một căn nhà trên sông nước.\\
\noindent B. Quảng cáo các mặt hàng điện tử được bán trên chợ nổi.\\
\noindent C. Phê phán lối sống xa hoa, lãng phí của người dân vùng sông nước.\\
\noindent D. Giải thích vì sao du khách thích mua đồ điện tử trên ghe.

\interlinepenalty=0


\vspace{0.3cm}
\interlinepenalty=10000

\noindent \textbf{Câu 24:} Nhận định nào đúng về hình thức treo ``bẹo'' được nêu ở đoạn (4)?\\
\noindent A. Ghe bán dựng một chậu cây bất kỳ để phân biệt với ghe mua, thay cho biển hiệu.\\
\noindent B. Mỗi loại hàng phải treo đúng loại cây tương ứng theo quy định của chợ.\\
\noindent C. Ghe mua dựng chậu cây để người bán nhận ra khách quen.\\
\noindent D. Chỉ được dùng vào những ngày chợ không hoạt động.

\interlinepenalty=0


\vspace{0.3cm}
\interlinepenalty=10000

\noindent \textbf{Câu 25:} Nhận định nào sau đây \textbf{không} phù hợp với nội dung văn bản?\\
\noindent A. Chợ nổi vẫn giữ nguyên vai trò chính yếu là nơi giao thương, không thay đổi theo thời gian.\\
\noindent B. Chợ nổi là một phần bản sắc văn hóa của vùng đồng bằng sông nước.\\
\noindent C. Vai trò ban đầu của chợ nổi gắn liền với giao thông đường thủy.\\
\noindent D. Hoạt động du lịch góp phần giúp chợ nổi tồn tại theo một hướng mới.

\interlinepenalty=0


\vspace{0.5cm}
% LECTURING-BLOCK: passage
\noindent\textbf{Dựa vào thông tin dưới đây để trả lời các câu từ 26 đến 30}% LECTURING-BLOCK-INSTRUCTION


\par\nobreak

(1) Một cuốn sách, bất kỳ nội dung nó ra làm sao, cứ để yên không cấm thì người ta coi như thường. Dầu có lắm người khích thích vì nó chăng nữa, song cũng còn có lắm người coi như thường. Đến cấm đi một cái, thì hết thảy ai nấy đều chú ý vào nó. Cấm đi, là muốn cho người ta đừng đọc, mà không ngờ lại làm cho người ta càng đọc!

(2) Người ta nói rằng một pho tượng gỗ đã sơn thếp rồi mà chưa khai quang điểm nhãn thì nó chưa linh, đến chừng khai quang điểm nhãn rồi thì nó mới linh. Ấy vậy, cấm cuốn sách nào, tức là khai quang điểm nhãn cho cuốn sách ấy.

(3) Cho nên ở ta đây đã có nhiều sách linh lắm, có nhiều cuốn đã hiển thánh rồi. \textit{Việt Nam vong quốc sử} và \textit{Hải ngoại huyết thơ} đã hiển thánh hai mươi năm nay; còn mới đây \textit{Một bầu tâm sự} cũng đòi đạp đồng ngang lên nữa.

(4) Thiên hạ họ mua sách cấm mắc tiền lắm, thường giá mắc gấp đôi lúc chưa cấm, lâu ngày, rồi đến gấp mười gấp trăm. Tức như dân An Nam là dân không ham đọc sách, hay tiếc tiền, mà đến sách cấm thì cũng trằn trọc trằn xa mua cho được. Mua rồi đọc chùng đọc vụng, ai biết đâu mà bắt.

(5) Người Tàu đương đời Mãn Thanh, vào khoảng Khương Hy, Kiền Long, có nhiều bộ sách bị cấm, mà cấm chừng nào thì người ta lại càng đọc và càng lấy làm thích chừng nấy. Có người đời bấy giờ đã nói: ``Các cái vui ở thế gian không cái nào bằng trong đêm có tuyết, đóng cửa lại mà đọc sách cấm''.

(6) Ở đất nhà Nam ta, chỉ có mây mù mà không có tuyết, làm cho giảm mất cái thú vui trong khi đọc sách cấm. Song le, đã có cái thú khác thế vào. Có người đã phát minh ra được và nói rằng: ``Không có cái thú gì bằng cái thú trong lúc trời mưa, đóng cửa cho chặt lại, rồi......đọc sách cấm!''
\begin{flushright}\small\textit{(Phan Khôi, Cấm sách, sách cấm, Đông Pháp thời báo, Sài Gòn, số 763, ngày 1/9/1928; nguồn: vi.wikisource.org, truy cập 08/10/2026; có lược bớt)}\end{flushright}

\par\nobreak

\interlinepenalty=10000

\noindent \textbf{Câu 26:} Ở đoạn (2), việc dùng hình ảnh ``khai quang điểm nhãn'' để nói về cuốn sách bị cấm có ý nghĩa gì?\\
\noindent A. Khuyên người đọc nên tôn kính những cuốn sách bị cấm như tượng thờ.\\
\noindent B. Ca ngợi nghi lễ truyền thống của người Việt khi làm tượng gỗ.\\
\noindent C. Ngầm cho rằng chính lệnh cấm làm tăng sức hút, khiến cuốn sách bỗng trở nên ``linh'' trong mắt người đọc.\\
\noindent D. Chê trách những người làm nghề chế tác tượng gỗ.

\interlinepenalty=0


\vspace{0.3cm}
\interlinepenalty=10000

\noindent \textbf{Câu 27:} Cách nói ``sách linh'', ``hiển thánh'', ``đạp đồng ngang lên'' ở đoạn (3) có tác dụng chủ yếu nào?\\
\noindent A. Mượn ngôn ngữ tín ngưỡng để giễu nhại hiệu quả ngược của việc cấm sách, tạo giọng điệu mỉa mai, hóm hỉnh.\\
\noindent B. Phê phán những người mê tín dị đoan khi đọc sách.\\
\noindent C. Miêu tả không khí linh thiêng của các buổi cúng lễ ở địa phương.\\
\noindent D. Thể hiện niềm tin của tác giả vào sức mạnh siêu nhiên của sách.

\interlinepenalty=0


\vspace{0.3cm}
\interlinepenalty=10000

\noindent \textbf{Câu 28:} Việc nhắc đến chuyện ``người Tàu đương đời Mãn Thanh'' ở đoạn (5) nhằm mục đích gì?\\
\noindent A. Đưa dẫn chứng lịch sử để củng cố luận điểm: cấm sách chừng nào người ta càng đọc chừng nấy.\\
\noindent B. So sánh trình độ học vấn của người Tàu với người Việt.\\
\noindent C. Phản biện quan điểm cho rằng sách cấm đều là sách hay.\\
\noindent D. Giới thiệu một triều đại có nhiều bộ sách quý trong lịch sử.

\interlinepenalty=0


\vspace{0.3cm}
\interlinepenalty=10000

\noindent \textbf{Câu 29:} Qua toàn bộ văn bản, quan điểm của tác giả về việc cấm sách là gì?\\
\noindent A. Lệnh cấm không đạt được mục đích ngăn người ta đọc, trái lại còn khiến sách bị cấm được tìm đọc nhiều hơn.\\
\noindent B. Lệnh cấm chỉ có hiệu quả khi được áp dụng thống nhất ở cả ba kỳ.\\
\noindent C. Lệnh cấm nên thay bằng việc tăng giá bán những cuốn sách bị cho là xấu.\\
\noindent D. Lệnh cấm là cần thiết để bảo vệ người đọc khỏi những cuốn sách có hại.

\interlinepenalty=0


\vspace{0.3cm}
\interlinepenalty=10000

\noindent \textbf{Câu 30:} Nội dung nào sau đây \textbf{không} được đề cập trong văn bản?\\
\noindent A. Người xưa từng cho rằng đọc sách cấm trong đêm tuyết là một thú vui.\\
\noindent B. Người ta thường đọc sách cấm một cách lén lút.\\
\noindent C. Sách bị cấm thường được bán với giá cao hơn lúc chưa cấm.\\
\noindent D. Nhà cầm quyền công bố rõ lý do cụ thể trước khi cấm từng cuốn sách.

\interlinepenalty=0

\vspace{0.5cm}
% LECTURING-PART: TA
\textbf{1.2. TIẾNG ANH} \textit{}

\par\nobreak

% LECTURING-BLOCK: vact_gap_fill
\noindent\textbf{Questions 31-35: Choose a suitable word or phrase (A, B, C, or D) to fill in each blank.}% LECTURING-BLOCK-INSTRUCTION


\par\nobreak

\interlinepenalty=10000

\noindent \textbf{Câu 31:} By the time the rescue team arrived, the hikers \underline{\hspace{1.5cm}} in the cave for nearly two days.\\
\noindent\makebox[0.48\textwidth][l]{A. have been trapped}\makebox[0.48\textwidth][l]{B. had been trapped}\\
\noindent\makebox[0.48\textwidth][l]{C. had trapped}\makebox[0.48\textwidth][l]{D. were trapping}

\interlinepenalty=0


\interlinepenalty=10000

\noindent \textbf{Câu 32:} The company decided to \underline{\hspace{1.5cm}} the launch of its new phone until next spring because of supply problems.\\
\noindent\makebox[0.24\textwidth][l]{A. put off}\makebox[0.24\textwidth][l]{B. put out}\makebox[0.24\textwidth][l]{C. put on}\makebox[0.24\textwidth][l]{D. put up}

\interlinepenalty=0


\interlinepenalty=10000

\noindent \textbf{Câu 33:} Not only \underline{\hspace{1.5cm}} the entrance exam, but she also won a scholarship to study abroad.\\
\noindent\makebox[0.24\textwidth][l]{A. passed she}\makebox[0.24\textwidth][l]{B. did she pass}\makebox[0.24\textwidth][l]{C. she did pass}\makebox[0.24\textwidth][l]{D. she passed}

\interlinepenalty=0


\interlinepenalty=10000

\noindent \textbf{Câu 34:} The \underline{\hspace{1.5cm}} of the new bridge took nearly three years because of the heavy rains.\\
\noindent\makebox[0.24\textwidth][l]{A. construction}\makebox[0.24\textwidth][l]{B. constructed}\makebox[0.24\textwidth][l]{C. construct}\makebox[0.24\textwidth][l]{D. constructive}

\interlinepenalty=0


\interlinepenalty=10000

\noindent \textbf{Câu 35:} She is not used to \underline{\hspace{1.5cm}} up early, so she often arrives late for her first class.\\
\noindent\makebox[0.24\textwidth][l]{A. get}\makebox[0.24\textwidth][l]{B. getting}\makebox[0.24\textwidth][l]{C. got}\makebox[0.24\textwidth][l]{D. have got}

\interlinepenalty=0



% LECTURING-BLOCK: vact_error_id
\noindent\textbf{Questions 36-40: Each of the following sentences has one error (A, B, C, or D). Find it.}% LECTURING-BLOCK-INSTRUCTION


\par\nobreak

\interlinepenalty=10000

\noindent \textbf{Câu 36:} The number of students \errorpart{A}{who enrolls} in online courses \errorpart{B}{has increased} \errorpart{C}{significantly} \errorpart{D}{over the past decade}.

\interlinepenalty=0


\interlinepenalty=10000

\noindent \textbf{Câu 37:} \errorpart{A}{Hardly had} the speaker \errorpart{B}{begun} to talk \errorpart{C}{than} the microphone \errorpart{D}{stopped working}.

\interlinepenalty=0


\interlinepenalty=10000

\noindent \textbf{Câu 38:} Each of the \errorpart{A}{volunteers} \errorpart{B}{were given} a certificate \errorpart{C}{at the end} of the \errorpart{D}{programme}.

\interlinepenalty=0


\interlinepenalty=10000

\noindent \textbf{Câu 39:} The manager, \errorpart{A}{along with} his assistants, \errorpart{B}{are planning} to attend \errorpart{C}{the conference} \errorpart{D}{next month}.

\interlinepenalty=0


\interlinepenalty=10000

\noindent \textbf{Câu 40:} The film, \errorpart{A}{which we watched} last night, was \errorpart{B}{so} \errorpart{C}{bored} that I \errorpart{D}{fell asleep}.

\interlinepenalty=0



% LECTURING-BLOCK: vact_restatement
\noindent\textbf{Questions 41-45: Which of the following best restates each of the given sentences?}% LECTURING-BLOCK-INSTRUCTION


\par\nobreak

\interlinepenalty=10000

\noindent \textbf{Câu 41:} It is believed that the new policy will reduce traffic congestion.\\
\noindent A. People believe traffic congestion will reduce the new policy.\\
\noindent B. The new policy is believed to reduce traffic congestion.\\
\noindent C. The new policy is believed to have reduced traffic congestion.\\
\noindent D. The new policy is believed reducing traffic congestion.

\interlinepenalty=0


\interlinepenalty=10000

\noindent \textbf{Câu 42:} She was so tired that she couldn't keep her eyes open.\\
\noindent A. She was so tired as to keep her eyes open.\\
\noindent B. She was too tired that she couldn't keep her eyes open.\\
\noindent C. She was tired enough to keep her eyes open.\\
\noindent D. She was too tired to keep her eyes open.

\interlinepenalty=0


\interlinepenalty=10000

\noindent \textbf{Câu 43:} Unless the fee is paid by Friday, your registration will be cancelled.\\
\noindent A. Provided the fee is not paid, your registration would have been cancelled by Friday.\\
\noindent B. If the fee is paid by Friday, your registration will be cancelled.\\
\noindent C. Your registration will be cancelled if the fee isn't paid by Friday.\\
\noindent D. Your registration won't be cancelled whether or not the fee is paid by Friday.

\interlinepenalty=0


\interlinepenalty=10000

\noindent \textbf{Câu 44:} I regret not telling her the truth earlier.\\
\noindent A. I wish I had told her the truth earlier.\\
\noindent B. I would rather I didn't tell her the truth earlier.\\
\noindent C. I wish I told her the truth earlier.\\
\noindent D. I hope I had told her the truth earlier.

\interlinepenalty=0


\interlinepenalty=10000

\noindent \textbf{Câu 45:} ``You should have checked the schedule before leaving,'' the teacher said to us.\\
\noindent A. The teacher suggested checking the schedule before we left.\\
\noindent B. The teacher denied that we had checked the schedule before leaving.\\
\noindent C. The teacher told us to check the schedule before leaving.\\
\noindent D. The teacher criticised us for not checking the schedule before leaving.

\interlinepenalty=0



% LECTURING-BLOCK: vact_passage_en
% LECTURING-BLOCK-TITLE: LIGHT POLLUTION
\noindent\textbf{Questions 46-52: Read the passage carefully.}% LECTURING-BLOCK-INSTRUCTION
\begin{center}\textbf{LIGHT POLLUTION}\end{center}% LECTURING-BLOCK-TITLE-LINE


\par\nobreak

(1) Light pollution is the alteration of the natural patterns of light and dark in ecosystems, resulting from the unnecessary and inefficient use of artificial lighting. Specific categories of light pollution include light trespass, over-illumination, glare, light clutter, and skyglow. A single offending light source often falls into more than one of these categories.

(2) Although this type of pollution can exist throughout the day, its effects are magnified during the night with the contrast of the sky's darkness. It has been estimated that 83\% of the world's people live under light-polluted skies and that 23\% of the world's land area is affected by skyglow.

(3) The area affected by artificial illumination continues to increase. A major side effect of urbanization, light pollution is blamed for compromising health, disrupting ecosystems, and spoiling aesthetic environments. Studies show that urban areas are more at risk. Globally, it has increased by at least 49\% from 1992 to 2017, and between 2014 and 2022 by 16\%, with uneven development, e.g. France reducing through policy its light pollution by 33\% in that time period.

(4) Solutions to light pollution are often steps like adjusting light fixtures or using more appropriate light bulbs. Further remediation can be done with more efforts to educate the public in order to push legislative change. However, because it is a man-made phenomenon, addressing its impacts on humans and the environment has political, social, and economic considerations.

(5) Light trespass occurs when unwanted light enters one's property. A common light trespass problem occurs when a strong light enters the window of one's home from the outside, causing problems such as sleep deprivation. A number of cities in the U.S. have developed standards for outdoor lighting to protect the rights of their citizens against light trespass.

(6) DarkSky International was started to reduce the light going up into the sky which reduces the visibility of stars (see Skyglow below). This is any light that is emitted more than 90$^\circ$ above nadir. By limiting light at this 90$^\circ$ mark they have also reduced the light output in the 80--90$^\circ$ range which creates most of the light trespass issues.
\begin{flushright}\small\textit{(Adapted from Wikipedia, ``Light pollution'', en.wikipedia.org, CC BY-SA 4.0, accessed 8 October 2026; selected paragraphs, text unchanged)}\end{flushright}

\par\nobreak

\interlinepenalty=10000

\noindent Choose an option (A, B, C, or D) that best answers each question.% LECTURING-BLOCK-INSTRUCTION

\noindent \textbf{Câu 46:} What is the passage mainly about?\\
\noindent A. The laws on outdoor lighting in the United States\\
\noindent B. The reasons why city dwellers can no longer see the stars\\
\noindent C. The nature, extent and possible remedies of unnecessary artificial light\\
\noindent D. The history of an organisation that protects the night sky

\interlinepenalty=0


\interlinepenalty=10000

\noindent \textbf{Câu 47:} According to paragraph 1, which of the following is NOT named as a specific category of light pollution?\\
\noindent\makebox[0.48\textwidth][l]{A. light reflection}\makebox[0.48\textwidth][l]{B. skyglow}\\
\noindent\makebox[0.48\textwidth][l]{C. glare}\makebox[0.48\textwidth][l]{D. light clutter}

\interlinepenalty=0


\interlinepenalty=10000

\noindent \textbf{Câu 48:} The word \textbf{magnified} in paragraph 2 is closest in meaning to \underline{\hspace{1.2cm}}.\\
\noindent\makebox[0.24\textwidth][l]{A. intensified}\makebox[0.24\textwidth][l]{B. delayed}\makebox[0.24\textwidth][l]{C. hidden}\makebox[0.24\textwidth][l]{D. blurred}

\interlinepenalty=0


\interlinepenalty=10000

\noindent \textbf{Câu 49:} What can be inferred from paragraph 3 about France?\\
\noindent A. Its cities caused most of the global increase after 2014.\\
\noindent B. Its policy shows that light pollution can fall even while the global level rises.\\
\noindent C. It has the lowest level of light pollution in the world.\\
\noindent D. It was the first country to measure light pollution.

\interlinepenalty=0


\interlinepenalty=10000

\noindent \textbf{Câu 50:} The word \textbf{it} in paragraph 4 refers to \underline{\hspace{1.2cm}}.\\
\noindent\makebox[0.48\textwidth][l]{A. artificial lighting fixtures}\makebox[0.48\textwidth][l]{B. the public}\\
\noindent\makebox[0.48\textwidth][l]{C. legislative change}\makebox[0.48\textwidth][l]{D. light pollution}

\interlinepenalty=0


\interlinepenalty=10000

\noindent \textbf{Câu 51:} According to paragraph 5, which of the following is a typical problem of light trespass?\\
\noindent A. A strong outside light entering a window and disturbing people's sleep\\
\noindent B. Stars becoming harder to see from the countryside\\
\noindent C. Public buildings using too much electricity\\
\noindent D. Drivers being dazzled by street lamps

\interlinepenalty=0


\vspace{0.3cm}

\interlinepenalty=10000

\noindent \textbf{Câu 52:} It can be inferred from paragraph 6 that \underline{\hspace{1.2cm}}.\\
\noindent A. light trespass is the only kind of light pollution that matters\\
\noindent B. DarkSky International only works on indoor lighting\\
\noindent C. most light trespass comes from light shining straight up\\
\noindent D. limiting light sent upwards also helped to reduce light trespass

\interlinepenalty=0



% LECTURING-BLOCK: vact_passage_en
% LECTURING-BLOCK-TITLE: THE URBAN HEAT ISLAND
\noindent\textbf{Questions 53-60: Read the passage carefully.}% LECTURING-BLOCK-INSTRUCTION
\begin{center}\textbf{THE URBAN HEAT ISLAND}\end{center}% LECTURING-BLOCK-TITLE-LINE


\par\nobreak

(1) The urban heat island (UHI) effect is a meteorological and climatological phenomenon in which urban areas experience significantly warmer temperatures than surrounding rural areas. The temperature difference is usually larger at night than during the day, and is most apparent when winds are weak, under block conditions, noticeably during the summer and winter.

(2) The main cause of the UHI effect is from the modification of land surfaces, while waste heat generated by energy usage is a secondary contributor. Urban areas occupy about 0.5\% of the Earth's land surface but host more than half of the world's population. As a population center grows, it tends to expand its area and increase its average temperature. Another key term is heat island; the term can be used to refer to any area that is relatively hotter than the surroundings, but generally refers to human-disturbed areas.

(3) Monthly rainfall is greater downwind of cities, partially due to the UHI. Increases in heat within urban centers increase the length of growing seasons, decrease air quality by increasing the production of pollutants such as ozone, and decrease water quality as warmer waters flow into area streams and put stress on their ecosystems.

(4) Not all cities have a distinct urban heat island, and the heat island characteristics depend strongly on the background climate of the area where the city is located. The impact in a city can significantly change based on its local environment. Heat can be reduced by tree cover and green space, which act as sources of shade and promote evaporative cooling. Other options include green roofs, passive daytime radiative cooling applications, ventilation corridors, the use of lighter-colored surfaces, and less absorptive building materials. These reflect more sunlight and absorb less heat.

(5) Climate change is not the cause of urban heat islands, but it is causing more frequent and more intense heat waves, which in turn amplify the urban heat island effect in cities. Compact and dense urban development may also increase the urban heat island effect, leading to higher temperatures and increased exposure.

(6) A definition of urban heat island is: ``The relative warmth of a city compared with surrounding rural areas.'' This relative warmth is caused by ``heat trapping due to land use, the configuration and design of the built environment, including street layout and building size, the heat-absorbing properties of urban building materials, reduced ventilation, reduced greenery and water features, and domestic and industrial heat emissions generated directly from human activities''.
\begin{flushright}\small\textit{(Adapted from Wikipedia, ``Urban heat island'', en.wikipedia.org, CC BY-SA 4.0, accessed 8 October 2026; selected paragraphs, text unchanged)}\end{flushright}

\par\nobreak

\interlinepenalty=10000

\noindent Choose an option (A, B, C, or D) that best answers each question.% LECTURING-BLOCK-INSTRUCTION

\noindent \textbf{Câu 53:} The best title for the passage is \underline{\hspace{1.2cm}}.\\
\noindent A. The Effects of Rural Land Use on Local Weather\\
\noindent B. Why Cities Are Warmer Than Their Surroundings\\
\noindent C. Green Roofs: The Only Cure for Hot Cities\\
\noindent D. How Climate Change Created Urban Heat Islands

\interlinepenalty=0


\interlinepenalty=10000

\noindent \textbf{Câu 54:} The word \textbf{apparent} in paragraph 1 is closest in meaning to \underline{\hspace{1.2cm}}.\\
\noindent\makebox[0.24\textwidth][l]{A. harmful}\makebox[0.24\textwidth][l]{B. obvious}\makebox[0.24\textwidth][l]{C. doubtful}\makebox[0.24\textwidth][l]{D. unusual}

\interlinepenalty=0


\interlinepenalty=10000

\noindent \textbf{Câu 55:} According to paragraph 2, \underline{\hspace{1.2cm}}.\\
\noindent A. cities become cooler as they expand\\
\noindent B. waste heat from energy use adds less to the UHI effect than changes to land surfaces\\
\noindent C. the UHI effect appears only in cities with more than a million people\\
\noindent D. urban areas cover about half of the Earth's land surface

\interlinepenalty=0


\interlinepenalty=10000

\noindent \textbf{Câu 56:} The word \textbf{it} in paragraph 2 refers to \underline{\hspace{1.2cm}}.\\
\noindent\makebox[0.48\textwidth][l]{A. a heat island}\makebox[0.48\textwidth][l]{B. a population center}\\
\noindent\makebox[0.48\textwidth][l]{C. the world's population}\makebox[0.48\textwidth][l]{D. the Earth's land surface}

\interlinepenalty=0


\interlinepenalty=10000

\noindent \textbf{Câu 57:} Which of the following is NOT mentioned in paragraph 3 as an effect of the UHI?\\
\noindent A. Warmer water flowing into streams\\
\noindent B. Longer growing seasons\\
\noindent C. More rainfall on the downwind side of cities\\
\noindent D. Reduced levels of ozone pollution

\interlinepenalty=0


\interlinepenalty=10000

\noindent \textbf{Câu 58:} What does the writer imply by saying that not all cities have a distinct urban heat island?\\
\noindent A. The effect disappears when the wind is strong.\\
\noindent B. Cities near the sea never become warmer than the countryside.\\
\noindent C. The local climate and environment affect how strong the effect is.\\
\noindent D. Only the largest cities can produce the effect.

\interlinepenalty=0


\interlinepenalty=10000

\noindent \textbf{Câu 59:} Based on paragraphs 4 and 5, which of the following is LEAST likely to reduce the UHI effect?\\
\noindent A. Using lighter-coloured surfaces\\
\noindent B. Making urban development more compact and dense\\
\noindent C. Building green roofs\\
\noindent D. Planting more trees and keeping green space

\interlinepenalty=0


\interlinepenalty=10000

\noindent \textbf{Câu 60:} According to paragraph 5, \underline{\hspace{1.2cm}}.\\
\noindent A. urban heat islands are the main cause of heat waves\\
\noindent B. climate change does not create urban heat islands but makes heat waves that strengthen them\\
\noindent C. there is no connection between climate change and urban heat\\
\noindent D. climate change is the main cause of urban heat islands

\interlinepenalty=0

\vspace{0.5cm}
% LECTURING-PART: TO
\textbf{PHẦN 2. TOÁN HỌC} \textit{}

% LECTURING-BLOCK: single

\interlinepenalty=10000

\noindent \textbf{Câu 61:} Một mảnh vườn hình chữ nhật có chu vi $56$ m. Nếu tăng chiều dài thêm $3$ m và giảm chiều rộng đi $3$ m thì diện tích mảnh vườn giảm $15$ m$^2$. Chiều dài ban đầu của mảnh vườn là\\
\noindent A. $13$ m. \hfill B. $14$ m. \hfill C. $15$ m. \hfill D. $16$ m.

\interlinepenalty=0


\vspace{0.3cm}
\interlinepenalty=10000

\noindent \textbf{Câu 62:} Cho phương trình $x^2-2(m-1)x+m^2-3m=0$ ($m$ là tham số thực). Có bao nhiêu giá trị nguyên của $m$ thuộc đoạn $[-10;10]$ để phương trình đã cho có hai nghiệm thực phân biệt?\\
\noindent A. $10$. \hfill B. $11$. \hfill C. $12$. \hfill D. $21$.

\interlinepenalty=0


\vspace{0.3cm}
\interlinepenalty=10000

\noindent \textbf{Câu 63:} Một người gửi $200$ triệu đồng vào ngân hàng với lãi suất $6\%$/năm, theo hình thức lãi kép (tiền lãi mỗi năm được nhập vào vốn để tính lãi cho năm sau). Sau ít nhất bao nhiêu năm thì người đó có được số tiền (cả vốn lẫn lãi) lớn hơn $300$ triệu đồng?\\
\noindent A. $6$ năm. \hfill B. $7$ năm. \hfill C. $8$ năm. \hfill D. $9$ năm.

\interlinepenalty=0


\vspace{0.3cm}
\interlinepenalty=10000

\noindent \textbf{Câu 64:} Cho $a=\log_2 3$ và $b=\log_3 5$. Khi đó $\log_{15}24$ bằng\\
\noindent\makebox[0.48\textwidth][l]{A. $\dfrac{a+3}{a(b+1)}$.}\makebox[0.48\textwidth][l]{B. $\dfrac{a+3}{ab+1}$.}\\[0.25cm]
\noindent\makebox[0.48\textwidth][l]{C. $\dfrac{a+3}{a+b}$.}\makebox[0.48\textwidth][l]{D. $\dfrac{3a+1}{a(b+1)}$.}

\interlinepenalty=0


\vspace{0.3cm}
\interlinepenalty=10000

\noindent \textbf{Câu 65:} Từ các chữ số $0,1,2,3,4,5$ lập được bao nhiêu số tự nhiên có bốn chữ số đôi một khác nhau và chia hết cho $5$?\\
\noindent A. $96$. \hfill B. $108$. \hfill C. $120$. \hfill D. $132$.

\interlinepenalty=0


\vspace{0.3cm}
\interlinepenalty=10000

\noindent \textbf{Câu 66:} Số nghiệm của phương trình $2\sin^2x-3\sin x+1=0$ trên đoạn $[0;2\pi]$ là\\
\noindent A. $2$. \hfill B. $3$. \hfill C. $4$. \hfill D. $5$.

\interlinepenalty=0


\vspace{0.3cm}
\interlinepenalty=10000

\noindent \textbf{Câu 67:} Giá trị của $\displaystyle\lim_{x\to1}\dfrac{\sqrt{x+3}-2}{x^2-1}$ bằng\\
\noindent\makebox[0.24\textwidth][l]{A. $0$.}\makebox[0.24\textwidth][l]{B. $\dfrac14$.}\makebox[0.24\textwidth][l]{C. $\dfrac18$.}\makebox[0.24\textwidth][l]{D. $\dfrac1{16}$.}

\interlinepenalty=0


\vspace{0.3cm}
\interlinepenalty=10000

\noindent \textbf{Câu 68:} Có bao nhiêu tiếp tuyến của đồ thị hàm số $y=x^3-3x+1$ song song với đường thẳng $y=9x-5$?\\
\noindent A. $0$. \hfill B. $1$. \hfill C. $2$. \hfill D. $3$.

\interlinepenalty=0


\vspace{0.3cm}
\interlinepenalty=10000

\noindent \textbf{Câu 69:} Diện tích hình phẳng giới hạn bởi parabol $y=x^2$ và đường thẳng $y=2x$ bằng\\
\noindent\makebox[0.24\textwidth][l]{A. $2$.}\makebox[0.24\textwidth][l]{B. $4$.}\makebox[0.24\textwidth][l]{C. $\dfrac43$.}\makebox[0.24\textwidth][l]{D. $\dfrac83$.}

\interlinepenalty=0


\vspace{0.3cm}
\interlinepenalty=10000

\noindent \textbf{Câu 70:} Cho cấp số nhân $(u_n)$ có $u_1=2$ và $u_4=54$. Tổng $5$ số hạng đầu tiên của cấp số nhân đó bằng\\
\noindent\makebox[0.24\textwidth][l]{A. $121$.}\makebox[0.24\textwidth][l]{B. $188$.}\makebox[0.24\textwidth][l]{C. $242$.}\makebox[0.24\textwidth][l]{D. $324$.}

\interlinepenalty=0


\vspace{0.3cm}
\interlinepenalty=10000

\noindent \textbf{Câu 71:} Cho hình lập phương $ABCD.A'B'C'D'$. Gọi $\alpha$ là góc giữa đường thẳng $A'C$ và mặt phẳng $(ABCD)$. Khi đó\\
\noindent\makebox[0.48\textwidth][l]{A. $\tan\alpha=\dfrac12$.}\makebox[0.48\textwidth][l]{B. $\tan\alpha=\dfrac{\sqrt2}{2}$.}\\[0.25cm]
\noindent\makebox[0.48\textwidth][l]{C. $\tan\alpha=\sqrt2$.}\makebox[0.48\textwidth][l]{D. $\tan\alpha=1$.}

\interlinepenalty=0


\vspace{0.3cm}
\interlinepenalty=10000

\noindent \textbf{Câu 72:} Gieo đồng thời hai con xúc xắc cân đối và đồng chất. Xác suất để tổng số chấm xuất hiện trên hai con xúc xắc là một số nguyên tố bằng\\
\noindent\makebox[0.24\textwidth][l]{A. $\dfrac13$.}\makebox[0.24\textwidth][l]{B. $\dfrac49$.}\makebox[0.24\textwidth][l]{C. $\dfrac5{12}$.}\makebox[0.24\textwidth][l]{D. $\dfrac7{18}$.}

\interlinepenalty=0


\vspace{0.5cm}
% LECTURING-BLOCK: math_context
\noindent\textbf{Dựa vào thông tin dưới đây để trả lời các câu từ 73 đến 74}% LECTURING-BLOCK-INSTRUCTION


\par\nobreak

Một xưởng sản xuất hai loại sản phẩm $X$ và $Y$. Mỗi sản phẩm $X$ cần $2$ giờ máy và $1$ kg nguyên liệu; mỗi sản phẩm $Y$ cần $1$ giờ máy và $3$ kg nguyên liệu. Mỗi ngày xưởng có tối đa $100$ giờ máy và $150$ kg nguyên liệu. Lợi nhuận thu được là $30$ nghìn đồng trên mỗi sản phẩm $X$ và $40$ nghìn đồng trên mỗi sản phẩm $Y$. Giả sử xưởng bán hết mọi sản phẩm làm ra.

\par\nobreak

\interlinepenalty=10000

\noindent \textbf{Câu 73:} Lợi nhuận lớn nhất mà xưởng có thể thu được trong một ngày là\\
\noindent\makebox[0.48\textwidth][l]{A. $2\,500$ nghìn đồng.}\makebox[0.48\textwidth][l]{B. $2\,600$ nghìn đồng.}\\
\noindent\makebox[0.48\textwidth][l]{C. $2\,200$ nghìn đồng.}\makebox[0.48\textwidth][l]{D. $2\,000$ nghìn đồng.}

\interlinepenalty=0


\vspace{0.3cm}
\interlinepenalty=10000

\noindent \textbf{Câu 74:} Nếu lợi nhuận trên mỗi sản phẩm $Y$ giảm còn $10$ nghìn đồng (lợi nhuận mỗi sản phẩm $X$ không đổi) thì phương án sản xuất mỗi ngày cho lợi nhuận lớn nhất là\\
\noindent A. $50$ sản phẩm $X$ và $0$ sản phẩm $Y$.\\
\noindent B. $40$ sản phẩm $X$ và $20$ sản phẩm $Y$.\\
\noindent C. $0$ sản phẩm $X$ và $50$ sản phẩm $Y$.\\
\noindent D. $30$ sản phẩm $X$ và $40$ sản phẩm $Y$.

\interlinepenalty=0


\vspace{0.5cm}
% LECTURING-BLOCK: math_context
\noindent\textbf{Dựa vào thông tin dưới đây để trả lời các câu từ 75 đến 76}% LECTURING-BLOCK-INSTRUCTION


\par\nobreak

Một quả bóng được thả rơi từ độ cao $10$ m xuống mặt sàn phẳng. Mỗi lần chạm sàn, quả bóng nảy lên đến độ cao bằng $\dfrac35$ độ cao của lần rơi ngay trước đó. Bỏ qua mọi lực cản.

\par\nobreak

\interlinepenalty=10000

\noindent \textbf{Câu 75:} Độ cao quả bóng nảy lên sau lần chạm sàn thứ ba là\\
\noindent\makebox[0.24\textwidth][l]{A. $1{,}296$ m.}\makebox[0.24\textwidth][l]{B. $2{,}16$ m.}\makebox[0.24\textwidth][l]{C. $3{,}6$ m.}\makebox[0.24\textwidth][l]{D. $6$ m.}

\interlinepenalty=0


\vspace{0.3cm}
\interlinepenalty=10000

\noindent \textbf{Câu 76:} Tổng quãng đường quả bóng đã đi được (tính từ lúc thả) cho đến khi dừng hẳn là\\
\noindent A. $25$ m. \hfill B. $30$ m. \hfill C. $40$ m. \hfill D. $50$ m.

\interlinepenalty=0


\vspace{0.5cm}
% LECTURING-BLOCK: math_context
\noindent\textbf{Dựa vào thông tin dưới đây để trả lời các câu từ 77 đến 78}% LECTURING-BLOCK-INSTRUCTION


\par\nobreak

Trong không gian $Oxyz$, cho hai điểm $A(1;2;-1)$, $B(3;0;1)$ và mặt phẳng $(P)\colon x+y-z+2=0$.

\par\nobreak

\interlinepenalty=10000

\noindent \textbf{Câu 77:} Phương trình mặt phẳng trung trực của đoạn thẳng $AB$ là\\
\noindent\makebox[0.48\textwidth][l]{A. $x+y-z-3=0$.}\makebox[0.48\textwidth][l]{B. $x-y+z+1=0$.}\\
\noindent\makebox[0.48\textwidth][l]{C. $x-y+z-1=0$.}\makebox[0.48\textwidth][l]{D. $2x-2y+2z-1=0$.}

\interlinepenalty=0


\vspace{0.3cm}
\interlinepenalty=10000

\noindent \textbf{Câu 78:} Gọi $M$ là điểm thuộc mặt phẳng $(P)$. Giá trị nhỏ nhất của $MA+MB$ bằng\\
\noindent\makebox[0.24\textwidth][l]{A. $2\sqrt{11}$.}\makebox[0.24\textwidth][l]{B. $6$.}\makebox[0.24\textwidth][l]{C. $2\sqrt{10}$.}\makebox[0.24\textwidth][l]{D. $4\sqrt3$.}

\interlinepenalty=0


\vspace{0.5cm}
% LECTURING-BLOCK: math_context
\noindent\textbf{Dựa vào thông tin dưới đây để trả lời các câu từ 79 đến 81}% LECTURING-BLOCK-INSTRUCTION


\par\nobreak

Cho hàm số $y=x^3-3x^2+mx+1$ với $m$ là tham số thực.

\par\nobreak

\interlinepenalty=10000

\noindent \textbf{Câu 79:} Tập hợp tất cả các giá trị của $m$ để hàm số đồng biến trên $\mathbb{R}$ là\\
\noindent\makebox[0.24\textwidth][l]{A. $m>3$.}\makebox[0.24\textwidth][l]{B. $m\le3$.}\makebox[0.24\textwidth][l]{C. $m\ge1$.}\makebox[0.24\textwidth][l]{D. $m\ge3$.}

\interlinepenalty=0


\vspace{0.3cm}
\interlinepenalty=10000

\noindent \textbf{Câu 80:} Có bao nhiêu giá trị nguyên của $m$ thuộc đoạn $[-5;5]$ để hàm số đã cho có hai điểm cực trị?\\
\noindent A. $7$. \hfill B. $8$. \hfill C. $9$. \hfill D. $11$.

\interlinepenalty=0


\vspace{0.3cm}
\interlinepenalty=10000

\noindent \textbf{Câu 81:} Với $m=0$, tổng của giá trị lớn nhất và giá trị nhỏ nhất của hàm số trên đoạn $[-1;3]$ bằng\\
\noindent\makebox[0.24\textwidth][l]{A. $-3$.}\makebox[0.24\textwidth][l]{B. $-2$.}\makebox[0.24\textwidth][l]{C. $-1$.}\makebox[0.24\textwidth][l]{D. $1$.}

\interlinepenalty=0


\vspace{0.5cm}
% LECTURING-BLOCK: math_context
\noindent\textbf{Dựa vào thông tin dưới đây để trả lời các câu từ 82 đến 84}% LECTURING-BLOCK-INSTRUCTION


\par\nobreak

Một lớp có $40$ học sinh, trong đó có $24$ nam và $16$ nữ. Cuối học kì, lớp có $14$ học sinh đạt loại giỏi, gồm $10$ nam và $4$ nữ.

\par\nobreak

\interlinepenalty=10000

\noindent \textbf{Câu 82:} Chọn ngẫu nhiên một học sinh của lớp. Xác suất để học sinh đó là nam hoặc đạt loại giỏi là\\
\noindent\makebox[0.48\textwidth][l]{A. $\dfrac34$.}\makebox[0.48\textwidth][l]{B. $\dfrac45$.}\\[0.25cm]
\noindent\makebox[0.48\textwidth][l]{C. $\dfrac7{10}$.}\makebox[0.48\textwidth][l]{D. $\dfrac{17}{20}$.}

\interlinepenalty=0


\vspace{0.3cm}
\interlinepenalty=10000

\noindent \textbf{Câu 83:} Chọn ngẫu nhiên đồng thời $3$ học sinh của lớp. Xác suất để trong $3$ học sinh được chọn có đúng $1$ học sinh đạt loại giỏi là\\
\noindent\makebox[0.48\textwidth][l]{A. $\dfrac{7}{19}$.}\makebox[0.48\textwidth][l]{B. $\dfrac{21}{76}$.}\\[0.25cm]
\noindent\makebox[0.48\textwidth][l]{C. $\dfrac{35}{76}$.}\makebox[0.48\textwidth][l]{D. $\dfrac{45}{76}$.}

\interlinepenalty=0


\vspace{0.3cm}
\interlinepenalty=10000

\noindent \textbf{Câu 84:} Chọn ngẫu nhiên một học sinh của lớp. Biết học sinh được chọn đạt loại giỏi, xác suất để học sinh đó là nữ là\\
\noindent\makebox[0.24\textwidth][l]{A. $\dfrac14$.}\makebox[0.24\textwidth][l]{B. $\dfrac27$.}\makebox[0.24\textwidth][l]{C. $\dfrac57$.}\makebox[0.24\textwidth][l]{D. $\dfrac1{10}$.}

\interlinepenalty=0


\vspace{0.5cm}
% LECTURING-BLOCK: math_context
\noindent\textbf{Dựa vào thông tin dưới đây để trả lời các câu từ 85 đến 87}% LECTURING-BLOCK-INSTRUCTION


\par\nobreak

Trong mặt phẳng tọa độ $Oxy$, cho tam giác $ABC$ có $A(1;2)$, $B(5;2)$, $C(1;5)$ (xem hình vẽ).

\begin{center}
\begin{adjustbox}{max width=0.95\linewidth, max totalheight=0.38\textheight}
\begin{tikzpicture}[scale=1.0, >=stealth]
  \draw[->, thick] (-0.4,0) -- (6.6,0) node[right] {$x$};
  \draw[->, thick] (0,-0.4) -- (0,6.0) node[above] {$y$};
  \foreach \i in {1,...,6} {\draw (\i,0.07) -- (\i,-0.07) node[below, font=\small] {$\i$};}
  \foreach \j in {1,...,5} {\draw (0.07,\j) -- (-0.07,\j) node[left, font=\small] {$\j$};}

  \draw[thick] (1,2) -- (5,2) -- (1,5) -- cycle;
  \draw (1,2.3) -- (1.3,2.3) -- (1.3,2);
  \fill (1,2) circle (1.5pt) node[below left] {$A$};
  \fill (5,2) circle (1.5pt) node[below right] {$B$};
  \fill (1,5) circle (1.5pt) node[above left] {$C$};
    % captions last so nothing paints over them
    % captions last so nothing paints over them
    \node[below left, font=\small] at (0,0) {$O$};
\end{tikzpicture}
\end{adjustbox}
\end{center}

\par\nobreak

\interlinepenalty=10000

\noindent \textbf{Câu 85:} Bán kính đường tròn nội tiếp tam giác $ABC$ bằng\\
\noindent A. $1$. \hfill B. $1{,}5$. \hfill C. $2$. \hfill D. $2{,}5$.

\interlinepenalty=0


\vspace{0.3cm}
\interlinepenalty=10000

\noindent \textbf{Câu 86:} Phương trình đường tròn ngoại tiếp tam giác $ABC$ là\\
\noindent A. $(x-3)^2+\left(y-\dfrac72\right)^2=\dfrac52$.\\[0.25cm]
\noindent B. $(x-3)^2+(y-3)^2=\dfrac{25}4$.\\[0.25cm]
\noindent C. $(x-1)^2+(y-2)^2=\dfrac{25}4$.\\[0.25cm]
\noindent D. $(x-3)^2+\left(y-\dfrac72\right)^2=\dfrac{25}4$.

\interlinepenalty=0


\vspace{0.3cm}
\interlinepenalty=10000

\noindent \textbf{Câu 87:} Gọi $D(x_D;y_D)$ là chân đường phân giác trong của góc $A$ (điểm $D$ thuộc cạnh $BC$). Giá trị $x_D+y_D$ bằng\\
\noindent\makebox[0.24\textwidth][l]{A. $6$.}\makebox[0.24\textwidth][l]{B. $7$.}\makebox[0.24\textwidth][l]{C. $\dfrac{39}7$.}\makebox[0.24\textwidth][l]{D. $\dfrac{45}7$.}

\interlinepenalty=0


\vspace{0.5cm}
% LECTURING-BLOCK: math_context
\noindent\textbf{Dựa vào thông tin dưới đây để trả lời các câu từ 88 đến 90}% LECTURING-BLOCK-INSTRUCTION


\par\nobreak

Cho hình chóp $S.ABC$ có đáy $ABC$ là tam giác vuông tại $B$, $AB=a$, $BC=a\sqrt3$; cạnh bên $SA$ vuông góc với mặt phẳng đáy và $SA=2a$ (xem hình vẽ).

\begin{center}
\begin{adjustbox}{max width=0.95\linewidth, max totalheight=0.38\textheight}
\begin{tikzpicture}[scale=1.0, >=stealth]
  \coordinate (A) at (0,0);
  \coordinate (B) at (3,-0.8);
  \coordinate (C) at (5,1);
  \coordinate (S) at (0,4);
  \draw[thick] (S) -- (A) -- (B) -- cycle;
  \draw[thick] (S) -- (C);
  \draw[thick] (B) -- (C);
  \draw[thick, dashed] (A) -- (C);
  % góc vuông tại B
  \draw (2.662,-0.71) -- (2.92,-0.476) -- (3.26,-0.566);
  \fill (S) circle (1.5pt) node[above] {$S$};
  \fill (A) circle (1.5pt) node[left] {$A$};
  \fill (B) circle (1.5pt) node[below] {$B$};
  \fill (C) circle (1.5pt) node[right] {$C$};
    % captions last so nothing paints over them
    % captions last so nothing paints over them
    \node[left] at (0,2) {$2a$};
\node[below left] at (1.5,-0.4) {$a$};
\node[below right] at (4,0.1) {$a\sqrt3$};
\end{tikzpicture}
\end{adjustbox}
\end{center}

\par\nobreak

\interlinepenalty=10000

\noindent \textbf{Câu 88:} Thể tích khối chóp $S.ABC$ bằng\\
\noindent\makebox[0.48\textwidth][l]{A. $\dfrac{\sqrt3a^3}{3}$.}\makebox[0.48\textwidth][l]{B. $\dfrac{2\sqrt3a^3}{3}$.}\\[0.25cm]
\noindent\makebox[0.48\textwidth][l]{C. $\dfrac{\sqrt3a^3}{6}$.}\makebox[0.48\textwidth][l]{D. $\sqrt3a^3$.}

\interlinepenalty=0


\vspace{0.3cm}
\interlinepenalty=10000

\noindent \textbf{Câu 89:} Góc giữa đường thẳng $SC$ và mặt phẳng $(ABC)$ bằng\\
\noindent\makebox[0.24\textwidth][l]{A. $30^\circ$.}\makebox[0.24\textwidth][l]{B. $90^\circ$.}\makebox[0.24\textwidth][l]{C. $45^\circ$.}\makebox[0.24\textwidth][l]{D. $60^\circ$.}

\interlinepenalty=0


\vspace{0.3cm}
\interlinepenalty=10000

\noindent \textbf{Câu 90:} Khoảng cách từ điểm $A$ đến mặt phẳng $(SBC)$ bằng\\
\noindent\makebox[0.48\textwidth][l]{A. $\dfrac{2\sqrt5a}{5}$.}\makebox[0.48\textwidth][l]{B. $\dfrac{\sqrt5a}{5}$.}\\[0.25cm]
\noindent\makebox[0.48\textwidth][l]{C. $\dfrac{\sqrt3a}{2}$.}\makebox[0.48\textwidth][l]{D. $\dfrac{2a}{3}$.}

\interlinepenalty=0

\vspace{0.5cm}
% LECTURING-PART: KH
\textbf{PHẦN 3. TƯ DUY KHOA HỌC\\ 3.1. LOGIC, PHÂN TÍCH SỐ LIỆU} \textit{}

\par\nobreak

% LECTURING-BLOCK: logic_puzzle
\noindent\textbf{Dựa vào thông tin dưới đây để trả lời các câu từ 91 đến 93}% LECTURING-BLOCK-INSTRUCTION


\par\nobreak

Sáu gian hàng $P, Q, R, S, T, U$ được xếp thành một hàng thẳng ở một hội chợ sách, đánh số vị trí từ $1$ đến $6$ tính từ cổng vào (mỗi vị trí có đúng một gian hàng). Việc sắp xếp phải thỏa mãn các điều kiện sau:
\begin{itemize}
\item Gian $P$ nằm trước gian $Q$ (vị trí của $P$ nhỏ hơn vị trí của $Q$).
\item Hai gian $R$ và $S$ nằm liền kề nhau.
\item Gian $T$ không nằm ở vị trí số $1$ hoặc vị trí số $6$.
\item Gian $U$ nằm trước gian $T$.
\item Hai gian $P$ và $T$ không nằm liền kề nhau.
\end{itemize}

\par\nobreak

\interlinepenalty=10000

\noindent \textbf{Câu 91:} Gian $P$ có thể nằm ở bao nhiêu vị trí khác nhau?\\
\noindent A. $3$. \hfill B. $4$. \hfill C. $5$. \hfill D. $6$.

\interlinepenalty=0


\vspace{0.3cm}
\interlinepenalty=10000

\noindent \textbf{Câu 92:} Nếu gian $T$ nằm ở vị trí số $5$ thì điều nào sau đây \textbf{phải} đúng?\\
\noindent\makebox[0.48\textwidth][l]{A. Gian $U$ nằm ở vị trí số $1$.}\makebox[0.48\textwidth][l]{B. Gian $P$ nằm ở vị trí số $1$.}\\
\noindent\makebox[0.48\textwidth][l]{C. Gian $P$ nằm trước gian $U$.}\makebox[0.48\textwidth][l]{D. Gian $Q$ nằm ở vị trí số $6$.}

\interlinepenalty=0


\vspace{0.3cm}
\interlinepenalty=10000

\noindent \textbf{Câu 93:} Nếu gian $R$ nằm ở vị trí số $1$ thì gian nào nằm ở vị trí số $4$?\\
\noindent\makebox[0.24\textwidth][l]{A. Gian $P$.}\makebox[0.24\textwidth][l]{B. Gian $Q$.}\makebox[0.24\textwidth][l]{C. Gian $U$.}\makebox[0.24\textwidth][l]{D. Gian $T$.}

\interlinepenalty=0


\vspace{0.5cm}
% LECTURING-BLOCK: logic_puzzle
\noindent\textbf{Dựa vào thông tin dưới đây để trả lời các câu từ 94 đến 96}% LECTURING-BLOCK-INSTRUCTION


\par\nobreak

Một câu lạc bộ có tám học sinh: $F, G, H, J, K, L, M, N$. Câu lạc bộ cần chọn đúng bốn học sinh vào đội tham dự cuộc thi. Việc chọn đội phải thỏa mãn các điều kiện sau:
\begin{itemize}
\item $F$ và $G$ không được cùng được chọn.
\item Nếu $H$ được chọn thì $J$ cũng phải được chọn.
\item Ba học sinh $K$, $J$, $L$ không được đồng thời cùng được chọn.
\item $L$ và $M$ không được cùng được chọn.
\item Nếu $G$ được chọn thì $M$ cũng phải được chọn.
\item Ít nhất một trong hai học sinh $K$ và $N$ phải được chọn.
\end{itemize}

\par\nobreak

\interlinepenalty=10000

\noindent \textbf{Câu 94:} Nhóm bốn học sinh nào dưới đây là một đội hợp lệ?\\
\noindent\makebox[0.24\textwidth][l]{A. $J, K, L, M$.}\makebox[0.24\textwidth][l]{B. $G, K, L, N$.}\makebox[0.24\textwidth][l]{C. $H, J, L, N$.}\makebox[0.24\textwidth][l]{D. $F, H, K, N$.}

\interlinepenalty=0


\vspace{0.3cm}
\interlinepenalty=10000

\noindent \textbf{Câu 95:} Nếu cả $J$ và $M$ đều được chọn thì điều nào sau đây \textbf{có thể} đúng?\\
\noindent\makebox[0.48\textwidth][l]{A. $H$ và $K$ cùng được chọn.}\makebox[0.48\textwidth][l]{B. $L$ cũng được chọn.}\\
\noindent\makebox[0.48\textwidth][l]{C. Cả $K$ và $N$ đều không được chọn.}\makebox[0.48\textwidth][l]{D. $F$ và $G$ cùng được chọn.}

\interlinepenalty=0


\vspace{0.3cm}
\interlinepenalty=10000

\noindent \textbf{Câu 96:} Nếu cả $H$ và $K$ đều không được chọn thì có bao nhiêu đội hợp lệ khác nhau?\\
\noindent A. $2$. \hfill B. $3$. \hfill C. $4$. \hfill D. $5$.

\interlinepenalty=0


\vspace{0.5cm}
% LECTURING-BLOCK: data_set
\noindent\textbf{Dựa vào thông tin dưới đây để trả lời các câu từ 97 đến 99}% LECTURING-BLOCK-INSTRUCTION


\par\nobreak

Hai biểu đồ tròn dưới đây cho biết cơ cấu doanh thu theo ba mảng hoạt động (Sản xuất, Thương mại, Dịch vụ) của công ty $X$ trong hai năm $2020$ và $2024$. Tổng doanh thu của công ty là $400$ tỉ đồng năm $2020$ và $600$ tỉ đồng năm $2024$.

\begin{center}
\begin{adjustbox}{max width=0.95\linewidth, max totalheight=0.38\textheight}
\begin{tabular}{@{}c@{}}
\begin{tikzpicture}[scale=0.9, align=center]
  % Năm 2020
  \draw[thick, fill=gray!15] (0,0) -- (90:1.5) arc (90:-72:1.5) -- cycle;
  \draw[thick, fill=gray!55] (0,0) -- (-72:1.5) arc (-72:-198:1.5) -- cycle;
  \draw[thick, fill=white] (0,0) -- (-198:1.5) arc (-198:-270:1.5) -- cycle;




  % Năm 2024
  \draw[thick, fill=gray!15] (9,0) -- ++(90:1.5) arc (90:-18:1.5) -- cycle;
  \draw[thick, fill=gray!55] (9,0) -- ++(-18:1.5) arc (-18:-162:1.5) -- cycle;
  \draw[thick, fill=white] (9,0) -- ++(-162:1.5) arc (-162:-270:1.5) -- cycle;
    % captions last so nothing paints over them
    % captions last so nothing paints over them
    \node[font=\small, align=center, fill=white, inner sep=1.5pt] at (3.1,0.9) {Sản xuất\\45\%};
\node[font=\small, align=center, fill=white, inner sep=1.5pt] at (-0.2,-2.7) {Thương mại\\35\%};
\node[font=\small, align=center, fill=white, inner sep=1.5pt] at (-3.1,0.9) {Dịch vụ\\20\%};
\node[font=\small\bfseries, fill=white, inner sep=1.5pt] at (0,2.6) {Năm 2020 (tổng 400 tỉ đồng)};
\node[font=\small, align=center, fill=white, inner sep=1.5pt] at (12.1,0.9) {Sản xuất\\30\%};
\node[font=\small, align=center, fill=white, inner sep=1.5pt] at (9,-2.7) {Thương mại\\40\%};
\node[font=\small, align=center, fill=white, inner sep=1.5pt] at (5.9,0.9) {Dịch vụ\\30\%};
\node[font=\small\bfseries, fill=white, inner sep=1.5pt] at (9,2.6) {Năm 2024 (tổng 600 tỉ đồng)};
\end{tikzpicture}
\end{tabular}
\end{adjustbox}
\end{center}

Tốc độ tăng trưởng kép hằng năm (CAGR) của một đại lượng từ giá trị đầu $V_0$ đến giá trị cuối $V_n$ sau $n$ năm được tính theo công thức $\mathrm{CAGR}=\left(\dfrac{V_n}{V_0}\right)^{\frac1n}-1$.

\par\nobreak

\interlinepenalty=10000

\noindent \textbf{Câu 97:} Doanh thu của mảng Dịch vụ năm $2024$ tăng bao nhiêu phần trăm so với năm $2020$?\\
\noindent\makebox[0.24\textwidth][l]{A. $125\%$.}\makebox[0.24\textwidth][l]{B. $90\%$.}\makebox[0.24\textwidth][l]{C. $50\%$.}\makebox[0.24\textwidth][l]{D. $150\%$.}

\interlinepenalty=0


\vspace{0.3cm}
\interlinepenalty=10000

\noindent \textbf{Câu 98:} Nhận định nào sau đây đúng?\\
\noindent A. Doanh thu của mảng Sản xuất giảm vì tỉ trọng của mảng này giảm.\\
\noindent B. Doanh thu của mảng Sản xuất không thay đổi dù tỉ trọng của mảng này giảm.\\
\noindent C. Mảng Thương mại có tốc độ tăng doanh thu nhanh nhất.\\
\noindent D. Doanh thu của mảng Thương mại tăng gấp đôi.

\interlinepenalty=0


\vspace{0.3cm}
\interlinepenalty=10000

\noindent \textbf{Câu 99:} CAGR của tổng doanh thu công ty $X$ trong giai đoạn $2020$--$2024$ gần nhất với giá trị nào sau đây?\\
\noindent\makebox[0.24\textwidth][l]{A. $12{,}5\%$/năm.}\makebox[0.24\textwidth][l]{B. $9{,}2\%$/năm.}\makebox[0.24\textwidth][l]{C. $8{,}5\%$/năm.}\makebox[0.24\textwidth][l]{D. $10{,}7\%$/năm.}

\interlinepenalty=0


\vspace{0.5cm}
% LECTURING-BLOCK: data_set
\noindent\textbf{Dựa vào thông tin dưới đây để trả lời các câu từ 100 đến 102}% LECTURING-BLOCK-INSTRUCTION


\par\nobreak

Khảo sát thời gian tự học ở nhà trong một tuần (đơn vị: giờ) của $60$ học sinh lớp $12$ cho kết quả sau:

\begin{center}
\begin{adjustbox}{max width=\linewidth}
\begin{tabular}{|c|c|c|c|c|c|}
\hline
Thời gian (giờ) & $[0;2)$ & $[2;4)$ & $[4;6)$ & $[6;8)$ & $[8;10)$ \\ \hline
Số học sinh & $6$ & $14$ & $22$ & $12$ & $6$ \\ \hline
\end{tabular}
\end{adjustbox}
\end{center}

\par\nobreak

\interlinepenalty=10000

\noindent \textbf{Câu 100:} Nhóm chứa trung vị của mẫu số liệu ghép nhóm trên là\\
\noindent\makebox[0.24\textwidth][l]{A. $[2;4)$.}\makebox[0.24\textwidth][l]{B. $[6;8)$.}\makebox[0.24\textwidth][l]{C. $[8;10)$.}\makebox[0.24\textwidth][l]{D. $[4;6)$.}

\interlinepenalty=0


\vspace{0.3cm}
\interlinepenalty=10000

\noindent \textbf{Câu 101:} Trung vị của mẫu số liệu ghép nhóm trên (làm tròn đến hàng phần trăm) là\\
\noindent A. $4{,}55$. \hfill B. $4{,}91$. \hfill C. $5{,}00$. \hfill D. $5{,}18$.

\interlinepenalty=0


\vspace{0.3cm}
\interlinepenalty=10000

\noindent \textbf{Câu 102:} Gọi $\bar{x}$ và $s$ lần lượt là số trung bình và độ lệch chuẩn của mẫu số liệu ghép nhóm trên (lấy giá trị đại diện của mỗi nhóm là trung điểm của nhóm đó). Giá trị $\bar{x}+s$ gần nhất với số nào sau đây?\\
\noindent A. $2{,}71$. \hfill B. $7{,}15$. \hfill C. $9{,}37$. \hfill D. $9{,}86$.

\interlinepenalty=0


\vspace{0.5cm}
\par\nobreak

% LECTURING-SUBHEADING
\noindent\textbf{3.2. SUY LUẬN KHOA HỌC}

% LECTURING-BLOCK: science
\noindent\textbf{Dựa vào thông tin dưới đây để trả lời các câu từ 103 đến 105}% LECTURING-BLOCK-INSTRUCTION


\par\nobreak

Dung dịch đệm là dung dịch chứa một axit yếu $\mathrm{HA}$ và bazơ liên hợp $\mathrm{A^-}$ của nó với nồng độ đáng kể; pH của dung dịch đệm thay đổi rất ít khi thêm một lượng nhỏ axit hoặc bazơ mạnh. Với dung dịch đệm, có thể tính pH theo công thức
\begin{center}$\mathrm{pH}=\mathrm{p}K_a+\lg\dfrac{[\mathrm{A^-}]}{[\mathrm{HA}]}$\end{center}
Trong thí nghiệm, người ta pha bốn dung dịch đệm, mỗi dung dịch có thể tích $1$ lít, từ axit axetic $\mathrm{CH_3COOH}$ (axit yếu, $\mathrm{p}K_a=4{,}76$) và natri axetat $\mathrm{CH_3COONa}$ theo bảng sau (coi thể tích không đổi khi thêm các chất; $\lg 2\approx 0{,}30$).

\begin{center}
\begin{tabular}{|c|c|c|}
\hline
Dung dịch & Số mol $\mathrm{CH_3COOH}$ & Số mol $\mathrm{CH_3COONa}$ \\ \hline
$Y_1$ & $0{,}20$ & $0{,}20$ \\ \hline
$Y_2$ & $0{,}05$ & $0{,}20$ \\ \hline
$Y_3$ & $0{,}20$ & $0{,}05$ \\ \hline
$Y_4$ & $0{,}10$ & $0{,}10$ \\ \hline
\end{tabular}
\end{center}

\par\nobreak

\interlinepenalty=10000

\noindent \textbf{Câu 103:} Giá trị pH của dung dịch $Y_4$ là\\
\noindent A. $2{,}88$. \hfill B. $4{,}76$. \hfill C. $5{,}76$. \hfill D. $7{,}00$.

\interlinepenalty=0


\vspace{0.3cm}
\interlinepenalty=10000

\noindent \textbf{Câu 104:} Sắp xếp pH của bốn dung dịch theo thứ tự tăng dần, ta được\\
\noindent\makebox[0.48\textwidth][l]{A. $Y_3<Y_1=Y_4<Y_2$.}\makebox[0.48\textwidth][l]{B. $Y_2<Y_1=Y_4<Y_3$.}\\
\noindent\makebox[0.48\textwidth][l]{C. $Y_3<Y_4<Y_1<Y_2$.}\makebox[0.48\textwidth][l]{D. $Y_1<Y_4<Y_3<Y_2$.}

\interlinepenalty=0


\vspace{0.3cm}
\interlinepenalty=10000

\noindent \textbf{Câu 105:} Thêm $0{,}05$ mol HCl vào dung dịch $Y_1$ (coi thể tích không đổi, HCl phản ứng hoàn toàn với $\mathrm{CH_3COO^-}$). Giá trị pH của dung dịch thu được gần nhất với giá trị nào sau đây?\\
\noindent\makebox[0.24\textwidth][l]{A. $3{,}00$.}\makebox[0.24\textwidth][l]{B. $4{,}54$.}\makebox[0.24\textwidth][l]{C. $4{,}76$.}\makebox[0.24\textwidth][l]{D. $4{,}98$.}

\interlinepenalty=0


\vspace{0.5cm}
% LECTURING-BLOCK: science
\noindent\textbf{Dựa vào thông tin dưới đây để trả lời các câu từ 106 đến 108}% LECTURING-BLOCK-INSTRUCTION


\par\nobreak

Để truyền tải điện năng đi xa, người ta dùng máy biến áp nâng hiệu điện thế ở nhà máy trước khi đưa lên đường dây. Khi truyền công suất $P$ đi trên đường dây có điện trở tổng cộng $R$ dưới hiệu điện thế $U$ (đo ở đầu đường dây), hệ số công suất $\cos\varphi$, công suất hao phí tỏa nhiệt trên đường dây là
\begin{center}$\Delta P=\dfrac{R\,P^2}{U^2\cos^2\varphi}$\end{center}
và hiệu suất truyền tải là $H=\dfrac{P-\Delta P}{P}$.

Một nhà máy phát điện truyền đi công suất $P=2{,}0$ MW trên đường dây có điện trở tổng cộng $R=10\ \Omega$ với hệ số công suất $\cos\varphi=1$. Lúc đầu hiệu điện thế ở đầu đường dây là $U=50$ kV.

\par\nobreak

\interlinepenalty=10000

\noindent \textbf{Câu 106:} Công suất hao phí trên đường dây lúc đầu bằng\\
\noindent A. $8$ kW. \hfill B. $16$ kW. \hfill C. $32$ kW. \hfill D. $160$ kW.

\interlinepenalty=0


\vspace{0.3cm}
\interlinepenalty=10000

\noindent \textbf{Câu 107:} Nếu nâng hiệu điện thế ở đầu đường dây lên $U'=100$ kV (các đại lượng khác không đổi) thì hiệu suất truyền tải bằng\\
\noindent\makebox[0.24\textwidth][l]{A. $99{,}5\%$.}\makebox[0.24\textwidth][l]{B. $99{,}9\%$.}\makebox[0.24\textwidth][l]{C. $99{,}8\%$.}\makebox[0.24\textwidth][l]{D. $99{,}2\%$.}

\interlinepenalty=0


\vspace{0.3cm}
\interlinepenalty=10000

\noindent \textbf{Câu 108:} Cho các phát biểu sau:\\
(1) Với $P$, $R$, $\cos\varphi$ không đổi, công suất hao phí tỉ lệ nghịch với bình phương hiệu điện thế $U$.\\
(2) Dùng dây dẫn có tiết diện lớn hơn (điện trở $R$ nhỏ hơn) thì làm giảm công suất hao phí.\\
(3) Muốn giảm công suất hao phí đi $4$ lần mà giữ nguyên $P$, $R$, $\cos\varphi$ thì phải tăng $U$ lên $4$ lần.\\
(4) Công suất hao phí không phụ thuộc vào hệ số công suất $\cos\varphi$.\\
Số phát biểu đúng là\\
\noindent A. $1$. \hfill B. $2$. \hfill C. $3$. \hfill D. $4$.

\interlinepenalty=0


\vspace{0.5cm}
% LECTURING-BLOCK: science
\noindent\textbf{Dựa vào thông tin dưới đây để trả lời các câu từ 109 đến 111}% LECTURING-BLOCK-INSTRUCTION


\par\nobreak

Ở nhiều loài thực vật, màu hoa do một gene có hai allele quy định: allele $A$ (trội hoàn toàn) quy định hoa đỏ, allele $a$ (lặn) quy định hoa trắng. Một quần thể ngẫu phối đang ở trạng thái cân bằng di truyền (Hardy--Weinberg) khi tần số các kiểu gene thỏa mãn $p^2\,AA+2pq\,Aa+q^2\,aa=1$, trong đó $p$, $q$ lần lượt là tần số của allele $A$ và $a$, $p+q=1$.

Các nhà nghiên cứu khảo sát hai quần thể của loài này, mỗi quần thể đều đang ở trạng thái cân bằng di truyền, và thu được số liệu sau:

\begin{center}
\begin{tabular}{|c|c|c|}
\hline
Quần thể & Tổng số cây khảo sát & Số cây hoa trắng \\ \hline
I & $1\,000$ & $160$ \\ \hline
II & $500$ & $45$ \\ \hline
\end{tabular}
\end{center}

\par\nobreak

\interlinepenalty=10000

\noindent \textbf{Câu 109:} Tần số allele $a$ của quần thể I là\\
\noindent A. $0{,}16$. \hfill B. $0{,}4$. \hfill C. $0{,}6$. \hfill D. $0{,}84$.

\interlinepenalty=0


\vspace{0.3cm}
\interlinepenalty=10000

\noindent \textbf{Câu 110:} Trong số các cây hoa đỏ của quần thể I, tỉ lệ cây có kiểu gene dị hợp ($Aa$) gần nhất với giá trị nào sau đây?\\
\noindent\makebox[0.24\textwidth][l]{A. $60{,}0\%$.}\makebox[0.24\textwidth][l]{B. $48{,}0\%$.}\makebox[0.24\textwidth][l]{C. $57{,}1\%$.}\makebox[0.24\textwidth][l]{D. $84{,}0\%$.}

\interlinepenalty=0


\vspace{0.3cm}
\interlinepenalty=10000

\noindent \textbf{Câu 111:} Nếu gộp toàn bộ cây của hai quần thể thành một quần thể mới (chưa xảy ra ngẫu phối), tần số allele $a$ của quần thể mới gần nhất với giá trị nào sau đây?\\
\noindent\makebox[0.24\textwidth][l]{A. $0{,}300$.}\makebox[0.24\textwidth][l]{B. $0{,}350$.}\makebox[0.24\textwidth][l]{C. $0{,}367$.}\makebox[0.24\textwidth][l]{D. $0{,}400$.}

\interlinepenalty=0


\vspace{0.5cm}
% LECTURING-BLOCK: science
\noindent\textbf{Dựa vào thông tin dưới đây để trả lời các câu từ 112 đến 114}% LECTURING-BLOCK-INSTRUCTION


\par\nobreak

Tỉ suất gia tăng dân số tự nhiên (đơn vị: \%) bằng hiệu số giữa tỉ suất sinh thô và tỉ suất tử thô (đơn vị: số người trên $1\,000$ dân) chia cho $10$. Tốc độ tăng dân số trung bình mỗi năm trong một giai đoạn $n$ năm được tính theo công thức $r=\left(\dfrac{N_n}{N_0}\right)^{\frac1n}-1$, trong đó $N_0$ và $N_n$ là dân số đầu và cuối giai đoạn. Theo ``quy tắc $70$'', nếu dân số tăng đều với tốc độ $r\%$ mỗi năm thì sau khoảng $\dfrac{70}{r}$ năm dân số sẽ tăng gấp đôi.

Bảng dưới đây là số liệu (giả định) của quốc gia $Z$.

\begin{center}
\begin{adjustbox}{max width=\linewidth}
\begin{tabular}{|c|c|c|c|}
\hline
Năm & Dân số (triệu người) & Sinh thô (người/$1\,000$ dân) & Tử thô (người/$1\,000$ dân) \\ \hline
$2015$ & $40$ & $15$ & $7$ \\ \hline
$2025$ & $44$ & $12$ & $8$ \\ \hline
\end{tabular}
\end{adjustbox}
\end{center}

\par\nobreak

\interlinepenalty=10000

\noindent \textbf{Câu 112:} Tỉ suất gia tăng dân số tự nhiên của quốc gia $Z$ năm $2025$ là\\
\noindent\makebox[0.24\textwidth][l]{A. $4{,}0\%$.}\makebox[0.24\textwidth][l]{B. $0{,}4\%$.}\makebox[0.24\textwidth][l]{C. $0{,}8\%$.}\makebox[0.24\textwidth][l]{D. $1{,}2\%$.}

\interlinepenalty=0


\vspace{0.3cm}
\interlinepenalty=10000

\noindent \textbf{Câu 113:} Tốc độ tăng dân số trung bình mỗi năm của quốc gia $Z$ trong giai đoạn $2015$--$2025$ gần nhất với giá trị nào sau đây?\\
\noindent\makebox[0.24\textwidth][l]{A. $1{,}2\%$.}\makebox[0.24\textwidth][l]{B. $0{,}96\%$.}\makebox[0.24\textwidth][l]{C. $1{,}0\%$.}\makebox[0.24\textwidth][l]{D. $0{,}4\%$.}

\interlinepenalty=0


\vspace{0.3cm}
\interlinepenalty=10000

\noindent \textbf{Câu 114:} Nếu từ năm $2025$ trở đi dân số của quốc gia $Z$ chỉ tăng tự nhiên với tỉ suất không đổi như năm $2025$ thì theo quy tắc $70$, sau khoảng bao nhiêu năm dân số tăng gấp đôi?\\
\noindent A. $100$ năm. \hfill B. $125$ năm. \hfill C. $175$ năm. \hfill D. $250$ năm.

\interlinepenalty=0


\vspace{0.5cm}
% LECTURING-BLOCK: science
\noindent\textbf{Dựa vào thông tin dưới đây để trả lời các câu từ 115 đến 117}% LECTURING-BLOCK-INSTRUCTION


\par\nobreak

Cách mạng công nghiệp lần thứ nhất khởi đầu ở nước Anh vào nửa sau thế kỉ XVIII. Nước Anh khi đó đã có nhiều điều kiện thuận lợi: sau cách mạng tư sản thế kỉ XVII, chế độ chính trị ổn định và bảo hộ quyền sở hữu; việc buôn bán và các thuộc địa mang lại nguồn vốn lớn cùng thị trường tiêu thụ rộng; trữ lượng than đá và quặng sắt dồi dào; tại nông thôn, nhiều nông dân mất đất nên trở thành nguồn lao động cho các xưởng thủ công và nhà máy.

Ngành dệt là ngành đi đầu với các máy kéo sợi, máy dệt. Việc cải tiến máy hơi nước của James Watt làm cho nhà máy không còn phụ thuộc vào sức nước của sông suối và có thể đặt ở bất cứ nơi nào thuận tiện cho việc thuê nhân công và vận chuyển. Về sau, máy hơi nước được lắp trên tàu thủy và đầu máy xe lửa, đường sắt nối các thành phố, các mỏ than và hải cảng với nhau. Cùng với sự ra đời của nhà máy, xã hội hình thành hai giai cấp cơ bản mới là tư sản công nghiệp và vô sản công nghiệp, đồng thời quá trình đô thị hóa diễn ra mạnh mẽ.

\par\nobreak

\interlinepenalty=10000

\noindent \textbf{Câu 115:} Điều kiện nào sau đây \textbf{không} được nêu trong đoạn thông tin như một thuận lợi giúp nước Anh khởi đầu cách mạng công nghiệp?\\
\noindent A. Sự viện trợ kĩ thuật và vốn từ các nước châu Á.\\
\noindent B. Nguồn vốn tích lũy từ buôn bán và thuộc địa.\\
\noindent C. Trữ lượng than đá và quặng sắt dồi dào.\\
\noindent D. Nguồn lao động từ những nông dân mất đất.

\interlinepenalty=0


\vspace{0.3cm}
\interlinepenalty=10000

\noindent \textbf{Câu 116:} Vì sao việc cải tiến máy hơi nước của James Watt có ý nghĩa quan trọng đối với sự phát triển của nhà máy?\\
\noindent A. Vì nhà máy không còn phụ thuộc vào vị trí có sức nước nên có thể đặt ở nơi thuận tiện về nhân công và vận chuyển.\\
\noindent B. Vì máy hơi nước thay thế hoàn toàn lao động của con người trong ngành dệt.\\
\noindent C. Vì máy hơi nước chỉ được dùng để vận chuyển hàng hóa trên sông.\\
\noindent D. Vì máy hơi nước giúp nước Anh có thêm nhiều thuộc địa.

\interlinepenalty=0


\vspace{0.3cm}
\interlinepenalty=10000

\noindent \textbf{Câu 117:} Một hệ quả xã hội quan trọng của cách mạng công nghiệp lần thứ nhất ở Anh là gì?\\
\noindent A. Sự xóa bỏ hoàn toàn khoảng cách giàu nghèo trong xã hội.\\
\noindent B. Sự hình thành hai giai cấp cơ bản: tư sản công nghiệp và vô sản công nghiệp.\\
\noindent C. Sự suy giảm của quá trình đô thị hóa.\\
\noindent D. Sự thay thế giai cấp tư sản bằng giai cấp quý tộc phong kiến.

\interlinepenalty=0


\vspace{0.5cm}
% LECTURING-BLOCK: science
\noindent\textbf{Dựa vào thông tin dưới đây để trả lời các câu từ 118 đến 120}% LECTURING-BLOCK-INSTRUCTION


\par\nobreak

Chị Lan mở một quán cà phê nhỏ. Mỗi tháng quán có chi phí cố định (thuê mặt bằng, lương nhân viên, khấu hao thiết bị) là $36$ triệu đồng. Chi phí biến đổi (nguyên liệu, ly, ống hút) là $18\,000$ đồng cho mỗi ly bán ra. Giá bán là $30\,000$ đồng một ly.

Điểm hòa vốn là sản lượng $Q^*$ mà tại đó tổng doanh thu bằng tổng chi phí:
\begin{center}$Q^*=\dfrac{\text{Chi phí cố định}}{\text{Giá bán}-\text{Chi phí biến đổi mỗi đơn vị}}$\end{center}
Lợi nhuận bằng tổng doanh thu trừ tổng chi phí.

\par\nobreak

\interlinepenalty=10000

\noindent \textbf{Câu 118:} Số ly cà phê quán cần bán mỗi tháng để hòa vốn là\\
\noindent\makebox[0.24\textwidth][l]{A. $4\,500$ ly.}\makebox[0.24\textwidth][l]{B. $2\,000$ ly.}\makebox[0.24\textwidth][l]{C. $3\,000$ ly.}\makebox[0.24\textwidth][l]{D. $3\,600$ ly.}

\interlinepenalty=0


\vspace{0.3cm}
\interlinepenalty=10000

\noindent \textbf{Câu 119:} Để đạt lợi nhuận $24$ triệu đồng mỗi tháng, quán cần bán bao nhiêu ly?\\
\noindent\makebox[0.24\textwidth][l]{A. $5\,000$ ly.}\makebox[0.24\textwidth][l]{B. $7\,000$ ly.}\makebox[0.24\textwidth][l]{C. $3\,000$ ly.}\makebox[0.24\textwidth][l]{D. $2\,000$ ly.}

\interlinepenalty=0


\vspace{0.3cm}
\interlinepenalty=10000

\noindent \textbf{Câu 120:} Giả sử chi phí cố định tăng $20\%$ và chị Lan đồng thời tăng giá bán lên $32\,000$ đồng một ly (chi phí biến đổi không đổi). Số ly cần bán mỗi tháng để hòa vốn lúc này gần nhất với\\
\noindent\makebox[0.24\textwidth][l]{A. $3\,000$ ly.}\makebox[0.24\textwidth][l]{B. $3\,600$ ly.}\makebox[0.24\textwidth][l]{C. $3\,086$ ly.}\makebox[0.24\textwidth][l]{D. $2\,571$ ly.}

\interlinepenalty=0\vspace{0.6cm}
\begin{center}
{\bfseries -{}-{}-{}-{}-{}-}\ \textbf{HẾT}\ {\bfseries -{}-{}-{}-{}-{}-}
\end{center}
\vspace{0.15cm}
{\small\itshape
Thí sinh không được sử dụng tài liệu. Cán bộ coi thi không giải thích gì thêm.}

\end{document}
